\documentclass[11pt]{article}
\usepackage{tcee}
\addbibresource{3B15.bib}
\tcetema[Regímenes cambiarios, intervención y regulación]{3.B.15}{Análisis comparado de los distintos regímenes cambiarios. Intervención y regulación de los mercados de cambio.}
\tcefecha{11/11/2026}
\begin{document}
\tceetitulo

\section*{Introducción}
\minutos{2,5}

\begin{itemize}
\item \textbf{\emph{Enganche}}:
  \begin{itemize}
  \item «Cada uno sabe dónde le aprieta el zapato», dice el refrán.
  \item A \autor{Hillel el Sabio}, el gran maestro judío del cambio de era, le pidieron que explicara toda la Torá mientras su interlocutor se sostenía sobre un pie. Respondió: «Lo que no quieras para ti, no lo hagas a tu prójimo: esta es toda la Torá; el resto es comentario».
  \item \autor{Michael Klein}, economista de la Universidad de Tufts, propuso resumir igual la macroeconomía internacional: los Gobiernos se enfrentan al trilema de la economía abierta, y todo lo demás es comentario \parencite{Economist2016Trilema}.
  \item Esta exposición es ese comentario: qué régimen cambiario elige cada país dentro del trilema, cómo lo sostiene interviniendo en el mercado y cuándo pone trabas a los capitales. La pregunta que la recorre es sencilla: \emph{¿hay un régimen cambiario mejor que los demás?}
  \end{itemize}

\item \textbf{\emph{Relevancia}}:
  \begin{itemize}
  \item El tipo de cambio enlaza la economía con el resto del mundo. Es probablemente el precio más importante de una economía, porque influye en todos los demás \vertema{3.B.13}:
    \begin{itemize}
    \item mueve la \emph{inflación}, por los precios de lo importado;
    \item mueve la \emph{competitividad} y el reparto de los recursos entre los bienes comerciables y los no comerciables;
    \item pone en juego la \emph{estabilidad financiera} cuando empresas, bancos o Estados deben en otra moneda.
    \end{itemize}
  \item Se fija en el mayor mercado del mundo: en abril de 2025 se negociaron 9,5 billones de dólares al día \parencite{BIS2026TrienalDatos}, lo que el comercio mundial de bienes y servicios mueve en menos de 4 días \parencite{OMC2026GTOS} \vertema{3.B.13}.
  \item El \concepto{régimen cambiario} decide quién fija ese precio y, con él, qué le queda a la política monetaria. Es una de las decisiones de política económica más importantes de un país abierto, y una de las que más crisis ha causado cuando se ha elegido o gestionado mal \vertema{3.B.17}.
  \end{itemize}

\item \textbf{\emph{Contextualización}}:
  \begin{itemize}
  \item \emph{Las definiciones de partida}:
    \begin{itemize}
    \item un \concepto{régimen cambiario} es el conjunto de reglas con que las autoridades determinan el valor de su moneda frente a las demás;\footnote{Usamos el \emph{tipo de cambio directo}, $E$: unidades de moneda nacional por unidad de moneda extranjera, es decir, el precio de la divisa. Si $E$ sube, la moneda nacional se \emph{deprecia}; si baja, se \emph{aprecia}. Con tipo fijo se habla de \emph{devaluación} y \emph{revaluación} cuando la autoridad cambia la paridad. El tipo de cambio real es $q\equiv E\,P^*/P$, y su subida es una depreciación real (una ganancia de competitividad-precio) \vertema{3.B.13}. Conviene decirlo al empezar: en los mercados gana peso la cotización indirecta (unidades de moneda extranjera por unidad de moneda nacional), y no todos los miembros del tribunal usan la misma.}
    \item con \concepto{tipo de cambio fijo}, la autoridad monetaria anuncia una paridad y se compromete a defenderla comprando y vendiendo divisas a ese precio, aunque puede devaluar o revaluar si lo considera necesario;
    \item con \concepto{tipo de cambio flexible}, la autoridad deja que el precio lo fijen la oferta y la demanda en el mercado de divisas;
    \item entre los 2 polos hay un continuo de \concepto{regímenes intermedios}, y lo que un país dice que hace (régimen \emph{de jure}) no siempre coincide con lo que hace (\emph{de facto}). Por eso el FMI clasifica los regímenes desde 1999 por lo que observa, en su informe anual sobre regímenes y restricciones cambiarias (AREAER), que publica desde 1950 \parencite{FMI2024AREAER2023}.
    \end{itemize}
  \item \emph{La historia, en 3 etapas}:
    \begin{itemize}
    \item el \emph{patrón oro}, hasta 1914 y en su restauración fallida de entreguerras: tipos fijos entre monedas convertibles en oro, dentro de los puntos del oro;
    \item \emph{Bretton Woods} (1944-1973): tipos fijos pero ajustables frente al dólar, con bandas del $\pm1$~\% y controles de capitales; el Acuerdo Smithsonian de diciembre de 1971 amplió las bandas al $\pm2{,}25$~\%, hasta que en marzo de 1973 las grandes monedas pasaron a flotar \parencite{CEA1972ERP} \vertema{3.B.20};
    \item desde 1973, el \concepto{«no sistema»}, como lo llamó \textcite{Williamson1976Benefits}: cada país elige su régimen y solo tiene que notificarlo al FMI, como quedó en la Segunda Enmienda de su Convenio Constitutivo, en vigor desde 1978. Muchos países dicen flotar pero intervienen (la \emph{flotación sucia}) \vertema{3.B.21}.
    \end{itemize}
  \item \emph{Una revisión reciente}: desde la crisis de 2008 se han revisado varias ideas sobre los regímenes, la intervención y la regulación.
    \begin{itemize}
    \item El FMI acepta en ciertos casos los controles de capitales (2012, revisado en 2022) y la intervención en el mercado de divisas (su Marco de Política Integrado, de 2020) \parencite{FMI2012VisionInstitucional,FMI2022RevisionVI,FMI2020IPF}.
    \item \autor{Hélène Rey} discute si el trilema no es en realidad un dilema \parencite{Rey2013Dilema}, y \autor{Gita Gopinath}, entonces economista jefe del FMI, si la depreciación sigue ganando competitividad cuando el comercio se factura en dólares \parencite{Gopinath2020DCP}.
    \end{itemize}
  \item \emph{España ha recorrido casi todo el menú}: tipos de cambio múltiples hasta el Plan de Estabilización de 1959 \vertema{4.A.27}, la peseta en las bandas del Sistema Monetario Europeo y, desde 1999, el euro, que es fijo frente a sus socios y flota frente al resto del mundo \vertema{3.B.16}.
  \end{itemize}

\item \textbf{\emph{Problemática}}:
  \begin{itemize}
  \item ¿Qué gana y qué pierde un país al fijar su tipo de cambio, y qué al dejarlo flotar?
  \item ¿De qué depende la elección, y qué dicen los datos?
  \item ¿Qué regímenes hay entre los 2 polos, y hacia cuáles van los países?
  \item ¿Para qué y cómo intervienen los bancos centrales en el mercado de divisas? ¿Sirve de algo?
  \item ¿Cuándo está justificado limitar los movimientos de capitales?
  \end{itemize}

\item \textbf{\emph{Estructura}}:
  \begin{itemize}
  \item La exposición va de lo que se elige a cómo se sostiene: primero el régimen y después sus instrumentos.
    \begin{itemize}
    \item \concepto{I. Diseño del régimen cambiario} (\ref{sec:diseno}): los 2 regímenes polares, sus ventajas y desventajas y los determinantes de la elección, con la evidencia.
    \item \concepto{II. Implementación del régimen cambiario} (\ref{sec:implementacion}): el menú de regímenes que clasifica el FMI y hacia dónde van los países.
    \item \concepto{III. Intervención y regulación} (\ref{sec:intervencion-regulacion}): las 2 herramientas con que se sostiene un régimen, la intervención en el mercado de divisas y la regulación de los movimientos de capitales.
    \end{itemize}
  \end{itemize}
\end{itemize}

\begin{esquema}[unbreakable]
\textbf{3.B.15. Regímenes cambiarios. Intervención y regulación}\\[2pt]
Introducción\esqtiempo{2,5}\\
I. Diseño del régimen cambiario: efectos sobre las variables macroeconómicas\esqtiempo{10,5}\\
\hspace*{1.5em}1. Idea\\
\hspace*{1.5em}2. Ventajas y desventajas de los tipos de cambio fijos y flexibles\\
\hspace*{1.5em}3. ¿Qué régimen arroja un mejor resultado? Determinantes y evidencia\\
II. Implementación del régimen cambiario\esqtiempo{6}\\
\hspace*{1.5em}1. Menú de regímenes cambiarios: grado de fijación y publicidad\\
\hspace*{1.5em}2. ¿Hacia qué régimen cambiario tienden los países?\\
III. Intervención y regulación de los mercados de cambio\esqtiempo{10}\\
\hspace*{1.5em}1. Intervención\\
\hspace*{1.5em}2. Regulación: los controles de capitales\\
Conclusión\esqtiempo{1}
\end{esquema}

\begin{notaopositor}
Reparto del cante: introducción, \tiempo{2,5}; I, \tiempo{10,5} (I.1, \tiempo{0,5}; I.2, \tiempo{7,5}: políticas de demanda y trilema, \tiempo{3,5}, equilibrio interno, \tiempo{1,5}, equilibrio externo, \tiempo{1}, mercados financieros, \tiempo{1}, y economía política, \tiempo{0,5}; I.3, \tiempo{2,5}); II, \tiempo{6} (II.1, \tiempo{4,5}; II.2, \tiempo{1,5}); III, \tiempo{10} (III.1, \tiempo{4,5}; III.2, \tiempo{5,5}), y conclusión, \tiempo{1}. Los gráficos imprescindibles son el trilema (y su versión del dilema), el eje de los regímenes, la zona objetivo y el modelo intertemporal con controles. España vuelve como ejemplo en cada apartado.
\end{notaopositor}


% ==========================================================================
\section{Diseño del régimen cambiario:\cortetitulo efectos sobre las variables macroeconómicas}\label{sec:diseno}
\minutos{10,5}

\begin{itemize}
\item Empezamos por la elección de fondo, la que el título llama análisis comparado: fijar el tipo de cambio o dejarlo flotar.
\item \emph{Conclusión anticipada}: ningún régimen es superior a los demás. El tipo fijo da credibilidad y estabilidad al comercio, pero deja el ajuste a precios y salarios y expone a ataques; el flexible absorbe las perturbaciones reales y deja la política monetaria libre, pero trae volatilidad. La elección depende de las características de cada país.
\end{itemize}

% --------------------------------------------------------------------------
\subsection{Idea}\label{sec:idea}

\begin{itemize}
\item El debate económico empezó por las ventajas y los inconvenientes de los 2 regímenes puros. Por eso comparamos 2 regímenes estilizados: un tipo fijo creíble y una flotación libre.
\item Los comparamos con \emph{libertad de movimientos de capitales}, porque hoy es difícil controlarlos durante mucho tiempo; lo que cambia con controles se ve al estudiar la regulación (\ref{sec:regulacion}).
\item Miramos el efecto del régimen en 5 dimensiones:
  \begin{itemize}
  \item la \emph{eficacia de las políticas de demanda} (\ref{sec:politicas-trilema});
  \item el \emph{equilibrio interno}: el ajuste de la renta y el empleo ante las perturbaciones, la inflación y la credibilidad (\ref{sec:equilibrio-interno});
  \item el \emph{equilibrio externo}: el comercio y la inversión, el riesgo de crisis y las reservas (\ref{sec:equilibrio-externo});
  \item el \emph{desarrollo de los mercados financieros} (\ref{sec:mercados-financieros});
  \item la \emph{economía política} (\ref{sec:economia-politica}).
  \end{itemize}
\item No sale un ganador: sale una lista de \emph{trade-offs}, que la tabla~\ref{tab:fijo-flexible} resume al final.
\end{itemize}

% --------------------------------------------------------------------------
\subsection{Ventajas y desventajas de los tipos de cambio fijos y flexibles}\label{sec:ventajas}

\subsubsection{Eficacia de las políticas de demanda y el trilema}\label{sec:politicas-trilema}

\begin{itemize}
\item El modelo de \autor{Mundell}-\autor{Fleming} da la respuesta de partida \parencite{Fleming1962Domestic,Mundell1963Capital} \vertema{3.A.32}.
  \begin{itemize}
  \item Añade al IS-LM una tercera curva, la BP (equilibrio de la balanza de pagos), con precios fijos y el tipo de cambio esperado constante.
  \item La \concepto{movilidad del capital} se mide con $\beta_K$, la sensibilidad de los flujos de capital al diferencial $i-i^*$:\footnote{Lleva subíndice porque $\beta$ es el factor de descuento del modelo intertemporal con controles de capitales (\ref{sec:regulacion}).} con movilidad nula, $\beta_K=0$, la BP es vertical; con movilidad perfecta, $\beta_K\to\infty$, es horizontal en $i=i^*$.
  \end{itemize}
\item Con movilidad perfecta, cada régimen deja viva una sola de las 2 políticas de demanda:
  \begin{itemize}
  \item \emph{con tipo fijo}, la política fiscal es muy eficaz y la monetaria, ineficaz. El banco central compra o vende todas las divisas que haga falta, y la cantidad de dinero se vuelve endógena: una expansión monetaria solo cambia la composición de su balance, con más crédito interno y menos reservas;
  \item \emph{con tipo flexible}, la política monetaria es muy eficaz, porque actúa a través de la depreciación, y la fiscal, ineficaz, porque la apreciación expulsa exportaciones netas.
  \end{itemize}
\item Con menos movilidad (de $\beta_K\to\infty$ a $\beta_K=0$) el resultado se matiza, pero el orden se mantiene (tabla~\ref{tab:mundell-fleming}):
  \begin{itemize}
  \item cuanto mayor es la movilidad, más eficaz es la política que funciona en cada régimen (la fiscal con tipo fijo, la monetaria con flexible) y menos la otra;
  \item con tipo fijo, la política monetaria no sirve con \emph{ninguna} movilidad: en ese régimen, el equilibrio interno queda supeditado al objetivo cambiario.
  \end{itemize}
\end{itemize}

\begin{table}[H]
\centering
\caption{Eficacia de las políticas monetaria y fiscal en el modelo de \autor{Mundell}-\autor{Fleming}}\label{tab:mundell-fleming}
\small
\begin{tabular}{>{\raggedright\arraybackslash}m{1.9cm}>{\raggedright\arraybackslash}m{3.0cm}>{\centering\arraybackslash}m{2.1cm}>{\centering\arraybackslash}m{2.1cm}>{\centering\arraybackslash}m{2.1cm}>{\centering\arraybackslash}m{2.1cm}}
\toprule
\rowcolor{tceeazulpalido} \multicolumn{2}{l}{Movilidad del capital} & Perfecta & Alta & Baja & Nula\\
\midrule
\multirow{2}{=}{Fiscal} & Tipo de cambio fijo & \textcolor{tceeverde}{Muy eficaz} & \textcolor{tceeverde}{Eficaz} & \textcolor{tceedorado}{Poco eficaz} & \textcolor{tceerojo}{Ineficaz}\\
 & Tipo de cambio flexible & \textcolor{tceerojo}{Ineficaz} & \textcolor{tceedorado}{Poco eficaz} & \textcolor{tceeverde}{Eficaz} & \textcolor{tceeverde}{Muy eficaz}\\
\midrule
\multirow{2}{=}{Monetaria} & Tipo de cambio fijo & \textcolor{tceerojo}{Ineficaz} & \textcolor{tceerojo}{Ineficaz} & \textcolor{tceerojo}{Ineficaz} & \textcolor{tceerojo}{Ineficaz}\\
 & Tipo de cambio flexible & \textcolor{tceeverde}{Muy eficaz} & \textcolor{tceeverde}{Eficaz} & \textcolor{tceeverde}{Eficaz} & \textcolor{tceeverde}{Eficaz}\\
\bottomrule
\end{tabular}
\fuente{Elaboración propia a partir de \textcite{Heijdra2017Foundations}. Movilidad alta: la BP es menos inclinada que la LM; baja: más inclinada. Con tipo fijo, la monetaria solo cambia las reservas, más deprisa cuanto mayor es la movilidad; con tipo flexible, su eficacia crece con la movilidad, y sin movilidad es la de la economía cerrada.}
\end{table}

\begin{notaopositor}
En el cante bastan las 2 políticas eficaces con movilidad perfecta: la fiscal con tipo fijo y la monetaria con tipo flexible. El resto de la tabla se cuenta de palabra.
\end{notaopositor}

\begin{itemize}
\item \emph{El precio de la eficacia monetaria con flotación es la volatilidad}: con precios rígidos a corto plazo y flexibles a largo, una expansión monetaria deprecia la moneda más a corto plazo que a largo, la \concepto{sobrerreacción} de \textcite{Dornbusch1976Expectations} \vertema{3.B.14}.
  \begin{itemize}
  \item Como los precios se mueven despacio, la volatilidad del tipo nominal pasa casi entera al real: con tipos flexibles, el tipo real es mucho más volátil que con tipos fijos \parencite{Mussa1986Nominal}, y mueve los precios relativos sin relación con los fundamentos.
  \item Para que esa volatilidad no tenga efectos reales grandes, el país necesita precios internos flexibles.
  \end{itemize}
\item Del resultado de \autor{Mundell}-\autor{Fleming} sale la regla de elección más famosa de la economía abierta.
\end{itemize}

\paragraph{El trilema de la economía abierta}\label{sec:trilema}

\subparagraph{a) Idea}

\begin{itemize}
\item El \concepto{trilema}\footnote{También se le llama \emph{trinidad imposible}, \emph{trinidad irreconciliable}, \emph{triángulo de la imposibilidad} o \emph{trilema monetario}.} convierte el resultado de \autor{Mundell}-\autor{Fleming} en una regla: de 3 objetivos deseables, un país solo puede tener 2 a la vez.
\item Se deduce del modelo de \textcite{Fleming1962Domestic} y \textcite{Mundell1963Capital}; el nombre de trilema se lo pusieron \autor{Maurice Obstfeld} y \autor{Alan Taylor} en un documento de trabajo de 1997, publicado en 1998 \parencite{ObstfeldTaylor1998Great,Economist2016Trilema}.
\item Ya lo intuyó \textcite{Keynes1930Treatise}: la libertad de los movimientos de capitales tensa cualquier orden monetario que quiera a la vez tipos de cambio estables y tipos de interés al servicio de la economía nacional.
\end{itemize}

\subparagraph{b) Modelo}

\begin{itemize}
\item \textbf{\emph{i) Supuestos}}:
  \begin{itemize}
  \item Un país pequeño, que toma como dado el tipo de interés mundial $i^*$.
  \item Activos nacionales y extranjeros sustitutos perfectos e inversores neutrales al riesgo: se cumple la paridad no cubierta de intereses (PNCI) con prima de riesgo nula, $\rho=0$ \vertema{3.B.13}.
  \end{itemize}
\item \textbf{\emph{ii) Desarrollo}}:
  \begin{itemize}
  \item Con libertad de capitales y un tipo fijo creíble, nadie espera que la moneda se deprecie ($\gamma E^e=0$), y la paridad de intereses obliga al banco central a seguir el tipo de interés del ancla:
    \funcionclave[eq:trilema]{i-i^*\approx\gamma E^e+\rho\quad\Rightarrow\quad i=i^*\ \ \text{si}\ \rho=0\ \text{y}\ \gamma E^e=0}
  \item De los 3 objetivos (un \emph{tipo de cambio fijo}, \emph{capitales que entran y salen sin trabas} y una \emph{política monetaria propia}), solo se pueden tener 2 a la vez; el tercero se cae \parencite{ObstfeldShambaughTaylor2005}. Pegatina: «no se puede nadar y guardar la ropa».\footnote{Es una aplicación de la regla de \textcite{Tinbergen1952Theory}: para alcanzar un número dado de objetivos hacen falta tantos instrumentos independientes como objetivos. Las autoridades querrían una política monetaria al servicio de los objetivos internos (la inflación, el paro), un tipo de cambio nominal estable y la plena movilidad del capital, que asigna bien el ahorro en el tiempo y reparte el riesgo; pero con libertad de capitales el tipo de interés es un solo instrumento y solo puede servir a un objetivo, el tipo de cambio o la economía propia. \autor{Jan Tinbergen} recibió en 1969 el primer Nobel de Economía, junto con \autor{Ragnar Frisch}.}
  \item Cada lado del triángulo es un régimen, que consigue los 2 objetivos de sus extremos y renuncia al del vértice opuesto (gráfico~\ref{fig:trilema}):
    \begin{itemize}
    \item \emph{tipo fijo y libertad de capitales}, sin política monetaria propia: la unión monetaria, la caja de conversión o la dolarización (Hong Kong, Dinamarca);
    \item \emph{política monetaria propia y libertad de capitales}, sin tipo fijo: la flotación (Estados Unidos, el euro frente al resto del mundo);
    \item \emph{tipo fijo y política monetaria propia}, sin libertad de capitales: Bretton Woods, con sus controles \vertema{3.B.20}.
    \end{itemize}
  \end{itemize}
\end{itemize}

\begin{figure}[H]
\centering
\caption{El trilema de la economía abierta}\label{fig:trilema}
\grafico[0.62\textwidth]{trilema}
\fuente{Elaboración propia a partir de \textcite{ObstfeldShambaughTaylor2005}. Cada vértice es un objetivo; cada lado, el régimen que consigue los 2 de sus extremos y renuncia al del vértice opuesto.}
\end{figure}

\begin{itemize}
\item \textbf{\emph{iii) Implicaciones}}:
  \begin{itemize}
  \item \emph{Un ejemplo}: si el banco central sube los tipos para frenar la inflación, entra capital y la moneda tiende a apreciarse. Le quedan 3 salidas:
    \begin{itemize}
    \item mantener la libertad de capitales y dejar que la moneda se aprecie: renuncia al tipo fijo;
    \item mantener el tipo fijo y limitar la entrada de capital: renuncia a la libertad de capitales;
    \item deshacer la subida de tipos: renuncia a la política monetaria propia.
    \end{itemize}
  \item Fijar el tipo de cambio con libertad de capitales es comprometerse a seguir la política monetaria del país ancla y, con ella, su ciclo. Si la producción diverge de la del ancla, el ajuste lo hacen los precios y los salarios o, si son rígidos, la producción y el empleo.
  \item Con tipo flexible, la producción y las variables nominales pueden separarse de las del ancla. Que eso sea bueno o malo depende de la credibilidad del banco central propio (\ref{sec:equilibrio-interno}).
  \item Volvamos a España: dentro del euro su tipo de cambio es fijo frente a los socios y el capital se mueve libremente, así que la política monetaria la decide el BCE para el conjunto. Es el trilema resuelto renunciando a la política monetaria propia \vertema{3.B.16}.
  \end{itemize}
\end{itemize}

\begin{notaopositor}
El trilema se demuestra con \autor{Mundell}-\autor{Fleming} (con tipo fijo y movilidad perfecta, la política monetaria es ineficaz) \vertema{3.A.32} o con el enfoque monetario con precios flexibles: con paridad del poder adquisitivo y paridad de intereses, fijar $E$ fija los precios relativos y el tipo de interés, $i=i^*$, y la cantidad de dinero pasa a ser endógena \vertema{3.B.14}.
\end{notaopositor}

\subparagraph{c) Extensiones}

\begin{itemize}
\item \emph{El trilema a medias}: las esquinas no son las únicas opciones. Con algo de flexibilidad cambiaria o con controles parciales, un país gana algo de autonomía monetaria; las esquinas del trilema están «redondeadas» \parencite{KleinShambaugh2015}.
\item \emph{El dilema}: la extensión más discutida pone en duda el tercer lado, que la flotación garantice la autonomía \parencite{Rey2013Dilema}. La vemos a continuación.
\end{itemize}

\subparagraph{d) Evidencia empírica}

\begin{itemize}
\item \emph{El trilema se ve en los datos}: con más de un siglo de datos, los países con tipo fijo y libre movilidad del capital siguen los tipos de interés del país de referencia; los que tienen controles de capitales o flotan conservan más autonomía monetaria \parencite{ObstfeldShambaughTaylor2005}.
\item Incluso con tipo fijo y capital libre queda algo de margen a corto plazo: los tipos de interés propios se separan un tiempo de los del ancla antes de seguirlos \parencite{KleinShambaugh2015}.
\end{itemize}

\subparagraph{e) Valoración}

\begin{itemize}
\item El trilema sigue siendo la referencia para pensar los regímenes: obliga a ver que la elección de un régimen es también la elección de una política monetaria.
\item Su punto débil está en el lado de la flotación: ¿basta con dejar flotar la moneda para tener una política monetaria propia?
\end{itemize}

\paragraph{¿Trilema o dilema? El ciclo financiero global}\label{sec:dilema}

\subparagraph{a) Idea}

\begin{itemize}
\item \textcite{Rey2013Dilema} pone en duda que la flotación garantice la autonomía: un país abierto a los flujos de capital y con tipo de cambio flexible no disfruta necesariamente de una autonomía monetaria plena.
\item La razón es un \concepto{ciclo financiero global}, común a todo el mundo, que no sigue las condiciones macroeconómicas de cada país \vertema{3.A.37} \vertema{3.A.42}.
\end{itemize}

\subparagraph{b) Modelo}

\begin{itemize}
\item \textbf{\emph{i) Supuestos}}:
  \begin{itemize}
  \item Libertad de movimientos de capitales.
  \item El dólar es la principal moneda de reserva y de financiación del mundo, y la política de la Reserva Federal decide cuánta liquidez en dólares hay.
  \item Los grandes intermediarios financieros se apalancan más cuando el riesgo medido es bajo y menos cuando sube.
  \end{itemize}
\item \textbf{\emph{ii) Desarrollo}}:
  \begin{itemize}
  \item Los precios de los activos con riesgo, los flujos brutos de capital y el crédito se mueven a la vez en todo el mundo, en sentido contrario al índice VIX, que mide la volatilidad esperada de la bolsa de Estados Unidos y sirve de termómetro de la aversión al riesgo \parencite{Rey2013Dilema}.
  \item Con modelos de vectores autorregresivos (VAR), uno de los principales determinantes de ese ciclo es la política monetaria del país central, la de la Reserva Federal \parencite{MirandaAgrippinoRey2020}.
  \item \emph{El mecanismo}: cuando la Reserva Federal baja los tipos, endeudarse en dólares se abarata, suben los precios de los activos en todo el mundo y, con ellos, el valor de las garantías con que se piden los préstamos. Las condiciones financieras se relajan en todas partes.\footnote{Según \textcite{Economist2016Trilema}, en uno de sus trabajos \autor{Rey} encuentra que una subida inesperada del tipo de interés de la Reserva Federal eleva pronto los diferenciales de los préstamos hipotecarios no solo en Estados Unidos, sino también en Canadá, el Reino Unido y Nueva Zelanda, países que dejan flotar su moneda.}
  \item El canal no es el tipo de interés de cada país, sino los precios de los activos, el apalancamiento de los bancos y la toma de riesgos. Pegatina: «cuando Estados Unidos estornuda, el mundo se resfría».
  \end{itemize}
\item \textbf{\emph{iii) Implicaciones}}:
  \begin{itemize}
  \item Con libertad de capitales, el ciclo financiero global limita la política monetaria nacional \emph{sea cual sea} el régimen cambiario.
  \item La elección real sería entre la libertad de capitales y una política monetaria propia: \emph{más que un trilema, un dilema} (gráfico~\ref{fig:dilema-rey}). La política monetaria propia solo sería plena si se controlan, directa o indirectamente, los movimientos de capitales.
  \item La política monetaria del país que flota sigue siendo eficaz, pero menos que si no hubiera libertad de capitales.
  \item Es un argumento a favor de las medidas macroprudenciales y de la coordinación internacional de las políticas \vertema{3.B.19}.
  \end{itemize}
\end{itemize}

\begin{figure}[H]
\centering
\caption{El dilema de \autor{Hélène Rey}}\label{fig:dilema-rey}
\grafico[0.62\textwidth]{dilema-rey}
\fuente{Elaboración propia a partir de \textcite{Rey2013Dilema}. Con libertad de capitales, el ciclo financiero global condiciona la política monetaria propia también con tipo flexible: el vértice del tipo fijo deja de ser la clave.}
\end{figure}

\subparagraph{c) Evidencia empírica}

\begin{itemize}
\item \emph{La evidencia matiza el dilema}:
  \begin{itemize}
  \item el ciclo financiero global se transmite más a los emergentes con tipo fijo que a los que flotan: el trilema se ha debilitado, pero sigue presente \parencite{ObstfeldOstryQureshi2019};
  \item la globalización financiera ha \emph{aumentado} la eficacia de la política monetaria de los que flotan, por los efectos de valoración de sus activos y pasivos en divisas \parencite{GeorgiadisMehl2016};\footnote{La eficacia de la política monetaria propia depende de 2 rasgos de los balances exteriores: su \emph{tamaño}, que mide cuánto llega el ciclo financiero global a las condiciones financieras internas y la debilita, y su \emph{exposición neta a las divisas}, que con tipo flexible crea efectos de valoración que la refuerzan. Dicho de otro modo, el ciclo financiero global puede atascar el canal del tipo de interés, pero refuerza el del tipo de cambio \parencite{GeorgiadisMehl2016}.}
  \item para la zona del euro, el BCE matiza que el motor principal es el apetito global por el riesgo y que no ha perdido el control de sus condiciones financieras \parencite{HabibVenditti2018}.
  \end{itemize}
\end{itemize}

\subparagraph{d) Valoración}

\begin{itemize}
\item El ciclo financiero global no hace indeseable la flotación; al contrario, la hace más valiosa. Si la Reserva Federal sube sus tipos y eso contrae la economía propia más de lo deseado, con tipo flexible el banco central puede bajar los suyos y dejar que la moneda se deprecie frente al dólar; con tipo fijo, no.
\item La flotación no aísla del todo, pero aísla más que el tipo fijo. Si no basta, la respuesta es complementarla con intervención, medidas macroprudenciales y medidas sobre los flujos (\ref{sec:intervencion-regulacion}), no volver a fijar.
\item La elección del régimen sigue siendo decisiva para la autonomía monetaria, también con un ciclo financiero global.
\end{itemize}

\begin{anotaciones}
\begin{itemize}
\item El argumento del dilema, contado por su autora: \parencite{Rey2013VoxEU}.
\item Un repaso del trilema y del ciclo financiero para no especialistas: \parencite{Economist2016Trilema}.
\end{itemize}
\end{anotaciones}

\subsubsection{Equilibrio interno}\label{sec:equilibrio-interno}

\begin{itemize}
\item Pasamos de las políticas de demanda al equilibrio interno: cómo se ajustan la renta y el empleo ante las perturbaciones y qué pasa con la inflación.
\item \emph{Ajuste ante perturbaciones reales}: es el argumento clásico de \textcite{Friedman1953Case} a favor de la flotación.
  \begin{itemize}
  \item Una perturbación real adversa (una caída de la demanda exterior o de la inversión) exige cambiar los precios relativos. Con precios y salarios flexibles, el ajuste llega por los precios y el régimen importa poco.
  \item Con precios rígidos, es más fácil mover un solo precio, el tipo de cambio, que miles: la depreciación abarata lo que se vende fuera y amortigua la caída de la producción.\footnote{\textcite{Friedman1953Case} lo comparaba con el horario de verano: es más fácil adelantar la hora que conseguir que todos adelanten a la vez sus costumbres.} Pegatina: «más vale doblarse que quebrarse».
  \item \emph{Un ejemplo}: ante una caída de la inversión que frena el crecimiento, la moneda tiende a depreciarse.
    \begin{itemize}
    \item \emph{Con tipo fijo}, el banco central tiene que subir los tipos para defender la paridad, lo que enfría más la inversión: la política monetaria no solo no estabiliza el ciclo, sino que puede agravarlo. Así les pasó a Argentina a finales de los 90 y a Francia en el Sistema Monetario Europeo en 1992-1993, que defendió el franco con tipos altos y no devaluó.
    \item \emph{Con tipo flexible}, el banco central no sube los tipos y deja que la moneda se deprecie, lo que evita que la inversión caiga más y sube las exportaciones: la caída de la producción se suaviza.
    \end{itemize}
  \item La defensa no siempre resiste: el 16 de septiembre de 1992, tras intervenciones masivas, el Reino Unido suspendió la participación de la libra en el mecanismo de cambios del Sistema Monetario Europeo \parencite{BoE1992Operation}; la peseta se devaluó varias veces entre 1992 y 1993 \parencite{EichengreenWyplosz1993Unstable} \vertema{3.B.17} \vertema{3.B.43}.
  \end{itemize}
\item \emph{Ajuste ante perturbaciones monetarias}: es al revés. Si cambia la demanda de dinero, con tipo fijo la cantidad de dinero se ajusta sola por las reservas y la renta no se mueve; con tipo flexible, la perturbación llega a la renta (\ref{sec:poole-abierta}).
\item \emph{Inflación y credibilidad}: con discrecionalidad, un banco central que quiere bajar el paro por debajo del natural acaba con más inflación y el mismo paro, el \concepto{sesgo inflacionista} de \textcite{BarroGordon1983Positive} \vertema{3.A.36}.
  \begin{itemize}
  \item \emph{Con tipo fijo}, la credibilidad se importa: fijar el tipo frente a la moneda de un banco central creíble baja el sesgo inflacionista sin subir el paro natural. Es «atarse las manos» \parencite{GiavazziPagano1988Advantage} \vertema{3.B.16}. Los planes de estabilización con ancla cambiaria de América Latina en los 90 (Argentina, Brasil, Chile) la usaron para bajar la inflación.
  \item \emph{El riesgo}: si la inflación propia converge despacio a la del ancla, el tipo real se aprecia,
    \begin{equation}\label{eq:apreciacion-real}
    \gamma_q\simeq\gamma_E+\pi^*-\pi=\pi^*-\pi<0\quad\text{si }\pi>\pi^*,
    \end{equation}
    y se pierde competitividad; recuperarla exige devaluar, que es justo lo que destruye la credibilidad que se buscaba \vertema{3.B.11}.
  \item \emph{Con tipo flexible}, la credibilidad tiene que nacer dentro: un banco central independiente y un objetivo de inflación \vertema{3.A.36}. Si falta, la inflación empuja a la depreciación, que encarece lo importado y alimenta más inflación; y con mucha inercia de precios aparece un efecto trinquete: la depreciación sube la inflación, pero la apreciación no la baja.
  \end{itemize}
\end{itemize}

\subsubsection{Equilibrio externo}\label{sec:equilibrio-externo}

\begin{itemize}
\item \emph{Comercio e inversión}: el tipo fijo elimina el riesgo de cambio y el coste de cubrirlo, y con ellos una incertidumbre que desanima las exportaciones, las importaciones y la inversión extranjera directa.
  \begin{itemize}
  \item \autor{John Maynard Keynes} defendía los tipos fijos por el comercio, pero criticó la vuelta del Reino Unido al patrón oro a la paridad de antes de la guerra: obligaba a los países deficitarios a subir tipos y bajar salarios, con más paro \parencite{Keynes1925Economic} \vertema{3.B.20}.
  \item Con una ecuación de gravedad, \textcite{Rose2000One} encontró que los países con la misma moneda comerciaban entre sí unas 3 veces más que con los demás; los metaanálisis rebajan mucho la cifra \parencite{RoseStanley2005Meta,GlickRose2016Currency} \vertema{3.B.10} \vertema{3.B.16}.\footnote{Se esperaba además que la unión monetaria transformara las estructuras productivas: sin la devaluación como salida, los Gobiernos tendrían más presión para hacer las reformas que antes habían rechazado \vertema{3.B.16}.}
  \item Los mercados financieros ofrecen hoy más instrumentos para cubrir el riesgo de cambio, y eso reduce la ventaja del tipo fijo.
  \end{itemize}
\item \emph{La ventaja del tipo flexible para el saldo exterior también se ha debilitado}. \autor{Friedman} y \autor{Mundell}-\autor{Fleming} suponen que la depreciación nominal se traslada casi uno a uno a la relación real de intercambio.
  \begin{itemize}
  \item \autor{Gita Gopinath}, entonces economista jefe del FMI, y sus coautores encuentran que una depreciación bilateral del 1~\% solo mueve un 0,1~\% la relación real de intercambio \parencite{BozGopinathPlagborgMoller2017}.
  \item La razón es que el comercio se factura en una moneda dominante, casi siempre el dólar, y no en la del exportador: los precios de exportaciones e importaciones son mucho más rígidos de lo que se creía y la depreciación mueve poco las exportaciones netas, sobre todo por el lado de las importaciones \parencite{Gopinath2020DCP} \vertema{3.B.14}.
  \item La flexibilidad sigue ayudando en los sectores que fijan sus precios en moneda local, como el turismo, y a los exportadores de materias primas cuando cae su precio.
  \end{itemize}
\item \emph{Riesgo de crisis}: un tipo fijo con capital libre ofrece al mercado una apuesta en un solo sentido. Hay 3 generaciones de modelos que lo explican \vertema{3.B.17}:
  \begin{itemize}
  \item en la \emph{primera}, una política incoherente con la paridad (un déficit financiado con crédito interno) agota las reservas hasta el ataque \parencite{Krugman1979Model};
  \item en la \emph{segunda}, el Gobierno abandona la paridad cuando defenderla le cuesta más que soltarla, y basta con que el mercado lo crea para que el ataque se produzca;
  \item en la \emph{tercera}, la devaluación hunde los balances de quienes deben en divisas y la crisis cambiaria se junta con la bancaria.
  \end{itemize}
\item \emph{Evidencia sobre las crisis}: en 50 emergentes entre 1980 y 2011, los regímenes intermedios son los más propensos a crisis bancarias y cambiarias; los anclajes duros no tanto, pero sí a los colapsos del crecimiento, quizá porque salir de ellos cuesta tanto que los Gobiernos esperan demasiado; y la variable que mejor anticipa la crisis es la sobrevaloración del tipo de cambio real \parencite{GhoshOstryQureshi2015}.
\item \emph{Reservas}: el tipo fijo exige muchas reservas para ser creíble, y mantenerlas tiene un coste de oportunidad. Con flotación hacen falta menos, aunque los emergentes que flotan también acumulan por precaución (\ref{sec:intervencion}). Cuántas bastan se mide con varios indicadores y con la métrica del FMI \vertema{3.B.22}.
\end{itemize}

\subsubsection{Desarrollo de los mercados financieros}\label{sec:mercados-financieros}

\begin{itemize}
\item Hasta aquí hemos comparado los regímenes por su relación con las políticas fiscal, monetaria y estructurales. Las crisis asiáticas de finales de los 90 añadieron otra dimensión: la crisis bancaria y la cambiaria se alimentaron una a otra.
\item \textcite{EichengreenHausmann1999} distinguen 3 explicaciones de por qué el régimen y la fragilidad financiera van juntos, y cada una pide un régimen distinto:
  \begin{itemize}
  \item \emph{Riesgo moral}: el tipo fijo es un seguro implícito. Bancos y empresas, que esperan un rescate si las cosas van mal, se endeudan en divisas y a corto plazo sin cubrirse, porque confían en la paridad. Pide más flexibilidad, para que todos aprendan a cubrirse, o la dolarización, que impide al banco central rescatar.
  \item \emph{Pecado original}: el país no puede endeudarse en su propia moneda a largo plazo, por su historial de inflación, por unos mercados poco desarrollados o por instituciones frágiles.\footnote{Endeudarse en dólares (o en euros en el este de Europa) es más barato que hacerlo en moneda local y, a veces, la única forma de conseguir financiación. Pero ata la política fiscal a la cambiaria: una devaluación puede llevar al impago de la deuda en divisas y a una crisis soberana, bancaria o las 2, según quién deba. Y basta con que suba el riesgo percibido de devaluación (o de redenominación, como en la crisis de la deuda soberana europea) para que el capital huya, el crédito se seque y el déficit público crezca.} Sus empresas acaban con \emph{descalces de moneda} (ingresos en moneda nacional y deudas en dólares) o \emph{de plazo} (proyectos largos financiados a corto) que no pueden cubrir: un país que solo puede deber en divisas no tiene con quién cubrirse \vertema{3.B.11}.
    \begin{itemize}
    \item \emph{Con flotación}, la depreciación hunde los balances de los deudores en dólares, y el miedo a que siga lleva a comprar divisas para cubrirse, lo que la agrava.
    \item \emph{Con tipo fijo o poco flexible}, defender la moneda exige subir los tipos, lo que lleva al impago a los deudores a corto plazo, frena la recuperación y, al hacer volátiles los tipos, impide que nazca un mercado de deuda local.
    \item Todas las salidas son imperfectas; los autores ven en la dolarización la menos mala.
    \end{itemize}
  \item \emph{Problemas de compromiso}: el que presta hoy cobra mañana, y lo que conviene a las 2 partes antes puede no convenir al deudor después; si las instituciones no garantizan el cobro, los contratos no se cumplen solos. La solución es institucional (mejores derechos de los acreedores) y el régimen importa menos: la flexibilidad deja al banco central sostener al sistema financiero en apuros, pero los acreedores, que lo anticipan, piden tipos más altos.
  \end{itemize}
\item \emph{Con flotación y credibilidad}, el mercado en moneda nacional se desarrolla: el sector público y el privado se endeudan más barato y se cubren en su propia moneda. Pero la flotación no evita todas las crisis: puede haber fugas de capitales o paradas súbitas aunque la moneda flote.
\item \emph{Hoy}: muchos emergentes se endeudan más en su moneda que hace 20 años, una cierta «redención» del pecado original. Pero el riesgo se ha desplazado a los inversores extranjeros en deuda local y a las empresas que emiten en dólares, y sigue alimentando el miedo a flotar (\ref{sec:determinantes}).
\end{itemize}

\begin{anotaciones}
\begin{itemize}
\item El trabajo original de las 3 hipótesis: \parencite{EichengreenHausmann1999}.
\end{itemize}
\end{anotaciones}

\subsubsection{Economía política}\label{sec:economia-politica}

\begin{itemize}
\item La última dimensión es política: el régimen reparte ganancias y pérdidas entre sectores y entre países \parencite{Frieden2015Currency}.
\item \emph{El tipo fijo traslada las tensiones dentro}: ante una perturbación, fuerza a bajar salarios y precios, y el ajuste es desigual.
  \begin{itemize}
  \item Lo soportan sobre todo los bienes no comerciables, porque el precio de los comerciables se fija en buena parte fuera (hoy, en dólares); los salarios de esos sectores bajan más.
  \item Con salarios nominales rígidos a la baja, el ajuste llega por las cantidades: más paro \parencite{SchmittGroheUribe2016Downward}.
  \item Fue la causa política de la caída del patrón oro de entreguerras: los ajustes de precios y salarios que exigía eran inviables con unos trabajadores y sindicatos cada vez más presentes en la vida política.
  \end{itemize}
\item \emph{La flotación traslada las tensiones fuera}: los países pueden caer en espirales de \concepto{devaluaciones competitivas}, un juego no cooperativo de suma negativa en el que cada uno intenta en vano ganar cuota a costa de sus socios (empobrecer al vecino) \vertema{3.B.19}.
  \begin{itemize}
  \item No es casualidad que al abandono del patrón oro en los años 30 le siguiera la época de las conferencias internacionales: fijar los tipos de cambio puede ser la salida cooperativa.
  \item Dentro del país, una depreciación reduce el poder adquisitivo de quienes consumen más bienes importados.
  \end{itemize}
\item Por eso el tipo fijo puede ser también una \emph{señal política} de integración o de cooperación con el país ancla, como lo fueron el Sistema Monetario Europeo y el euro.
\item Recapitulando, ninguno de los 2 regímenes gana en todas las dimensiones (tabla~\ref{tab:fijo-flexible}). En esta tesitura, ¿qué régimen debe elegir un país?
\end{itemize}

\begin{table}[H]
\centering
\caption{Tipo de cambio fijo frente a flexible}\label{tab:fijo-flexible}
\small
\begin{tabular}{>{\raggedright\arraybackslash}m{3.2cm}>{\raggedright\arraybackslash}m{5.6cm}>{\raggedright\arraybackslash}m{5.6cm}}
\toprule
\rowcolor{tceeazulpalido}\textbf{Dimensión} & \textbf{Tipo fijo} & \textbf{Tipo flexible}\\
\midrule
Políticas de demanda & \textcolor{tceeverde}{Fiscal eficaz}; \textcolor{tceerojo}{monetaria sin autonomía} & \textcolor{tceeverde}{Monetaria eficaz}; \textcolor{tceerojo}{fiscal poco eficaz}\\
Perturbaciones reales & \textcolor{tceerojo}{El ajuste recae en la renta y el empleo} & \textcolor{tceeverde}{Las absorbe el tipo de cambio}\\
Perturbaciones monetarias & \textcolor{tceeverde}{Las absorben las reservas} & \textcolor{tceerojo}{Llegan a la renta}\\
Inflación & \textcolor{tceeverde}{Credibilidad importada}; \textcolor{tceerojo}{riesgo de apreciación real} & Depende del banco central propio\\
Volatilidad, comercio e inversión & \textcolor{tceeverde}{Sin riesgo de cambio} & \textcolor{tceerojo}{Volatilidad y sobrerreacción}\\
Riesgo de crisis & \textcolor{tceerojo}{Ataques especulativos} & \textcolor{tceeverde}{Sin crisis cambiaria}, pero no sin paradas súbitas\\
Reservas & \textcolor{tceerojo}{Muchas, con su coste} & \textcolor{tceeverde}{Menos}\\
Sistema financiero & \textcolor{tceerojo}{Riesgo moral; descalces sin cubrir} & \textcolor{tceerojo}{Con deuda en dólares, efecto balance}\\
Economía política & Tensiones internas (ajuste de salarios); señal de integración & Tensiones externas (devaluaciones competitivas)\\
\bottomrule
\end{tabular}
\fuente{Elaboración propia a partir de \textcite{Friedman1953Case}, \textcite{Mundell1963Capital}, \textcite{GiavazziPagano1988Advantage}, \textcite{EichengreenHausmann1999} y \textcite{GhoshOstryQureshi2015}. En verde, ventaja; en rojo, inconveniente.}
\end{table}

% --------------------------------------------------------------------------
\subsection{¿Qué régimen arroja un mejor resultado?\cortetitulo Determinantes y evidencia}\label{sec:eleccion}

\begin{itemize}
\item Hemos visto qué cambia con cada régimen; ahora, de qué depende cuál conviene y qué dicen los datos.
\item \emph{Conclusión anticipada}: el régimen adecuado depende de dónde vienen las perturbaciones, de lo integrado que esté el país con su ancla, de la credibilidad de sus instituciones y de su desarrollo financiero. Los datos no dan un ganador universal. Pegatina, la del enganche: «cada uno sabe dónde le aprieta el zapato».
\end{itemize}

\subsubsection{Los determinantes de la elección}\label{sec:determinantes}

\begin{itemize}
\item \emph{La integración con el ancla}: fijar el tipo de cambio cuesta menos cuanto menos falta hace el tipo de cambio para ajustarse. Son los criterios de la teoría de las áreas monetarias óptimas \parencite{Mundell1961Theory,McKinnon1963Optimum,Kenen1969Theory} \vertema{3.B.16}:
  \begin{itemize}
  \item el \emph{tamaño y la apertura}: una economía pequeña y abierta gana más con la estabilidad del precio de lo que comercia y pierde menos, porque la depreciación se le va enseguida a los precios;
  \item las \emph{perturbaciones simétricas} con el ancla: si los ciclos van a la par, la política monetaria del ancla también le sirve;
  \item \emph{otros mecanismos de ajuste}: precios y salarios flexibles, movilidad del trabajo, transferencias fiscales;
  \item la \emph{diversificación} de la producción: las perturbaciones de cada sector se compensan.
  \end{itemize}
\item \emph{El origen de las perturbaciones}: reales o monetarias, como en el análisis de \autor{Poole} (\ref{sec:poole-abierta}).
\item \emph{La credibilidad}: un país con historial de inflación y un banco central poco creíble gana con un ancla dura; uno con instituciones creíbles puede permitirse flotar con un objetivo de inflación (\ref{sec:equilibrio-interno}).
\item \emph{El desarrollo financiero}: sin mercados de divisas profundos ni deuda en moneda propia, la flotación es volátil y la depreciación daña los balances. Por eso muchos emergentes que dicen flotar no se atreven a hacerlo: el \concepto{miedo a flotar} \parencite{CalvoReinhart2002}.
\item \emph{La economía política}: el régimen lo eligen Gobiernos sometidos a grupos de interés, como los sectores comerciables o los deudores en divisas, y pesa el ciclo electoral: las devaluaciones tienden a retrasarse hasta después de las elecciones \parencite{Frieden2015Currency}.
\item \emph{La movilidad del capital}: con controles, el trilema deja más margen para combinar un tipo estable con una política monetaria propia (\ref{sec:regulacion}).
\end{itemize}

\subsubsection{El origen de las perturbaciones:\cortetitulo \autor{Poole} en economía abierta}\label{sec:poole-abierta}

\paragraph{a) Idea}

\begin{itemize}
\item \textcite{Poole1970Optimal} mostró que un banco central debe fijar la cantidad de dinero o el tipo de interés según dónde nazcan las perturbaciones \vertema{3.A.36}.
\item En economía abierta, la elección es entre fijar el tipo de cambio o la cantidad de dinero, y la respuesta es parecida \parencite{Boyer1978Optimal}: con movilidad perfecta, fijar el tipo de cambio es fijar el tipo de interés en $i^*$.
\end{itemize}

\paragraph{b) Modelo}

\subparagraph{i) Supuestos}

\begin{itemize}
\item \autor{Mundell}-\autor{Fleming} con precios fijos y movilidad perfecta del capital, $i=i^*$ \vertema{3.A.32}.
\item 2 perturbaciones independientes, de media cero: una real, $\varepsilon^{\mathrm{IS}}$, en la demanda de bienes, con varianza $\sigma^2_{\mathrm{IS}}$, y una monetaria, $\varepsilon^{\mathrm{LM}}$, en la demanda de dinero, con varianza $\sigma^2_{\mathrm{LM}}$.
\item El objetivo es estabilizar la renta: elegir el régimen con menor $\mathrm{Var}(Y)$.
\end{itemize}

\subparagraph{ii) Desarrollo}

\begin{itemize}
\item \emph{Con tipo fijo}, el dinero es endógeno y las reservas absorben la perturbación monetaria. La real llega entera a la renta con el multiplicador de la economía abierta, porque ni el tipo de interés ni el tipo real ($\bar q=\overline E\,P^*/P$) se mueven:
  $$Y=\frac{\bar A-b\,i^*+x\,Y^*+\nu\,\bar q+\varepsilon^{\mathrm{IS}}}{1-c\,(1-t)+m}\quad\Rightarrow\quad \mathrm{Var}(Y)_{\text{fijo}}=\frac{\sigma^2_{\mathrm{IS}}}{[1-c\,(1-t)+m]^2}.$$
\item \emph{Con tipo flexible}, la cantidad de dinero es fija y la renta la decide la LM con $i=i^*$. La perturbación real la absorbe el tipo de cambio, que se aprecia o se deprecia hasta compensarla, pero la monetaria llega a la renta:
  $$\frac{M}{P}=k\,Y-h\,i^*+\varepsilon^{\mathrm{LM}}\quad\Rightarrow\quad Y=\frac{M/P+h\,i^*-\varepsilon^{\mathrm{LM}}}{k},\qquad \mathrm{Var}(Y)_{\text{flexible}}=\frac{\sigma^2_{\mathrm{LM}}}{k^2}.$$
\end{itemize}

\subparagraph{iii) Implicaciones}

\begin{itemize}
\item Conviene fijar el tipo de cambio si dominan las perturbaciones monetarias, y flotar si dominan las reales:
  \funcionclave[eq:poole-abierta]{\text{fijar }E\ \text{es mejor}\iff\frac{\sigma^2_{\mathrm{LM}}}{k^2}>\frac{\sigma^2_{\mathrm{IS}}}{[1-c\,(1-t)+m]^2}}
\item Es el argumento de \textcite{Friedman1953Case} visto desde el otro lado: las perturbaciones reales piden flexibilidad, y los países con una demanda de dinero inestable (inflación alta, innovación financiera, dolarización parcial) ganan con el tipo fijo.
\item Pegatina, la de \autor{Poole}: «perturbación real, fija la cantidad; monetaria, fija el tipo»; aquí, el tipo de cambio.
\end{itemize}

\paragraph{c) Extensiones}

\begin{itemize}
\item \emph{Regímenes intermedios}: si conviven las 2 perturbaciones, lo óptimo es una regla de intervención que acompaña solo en parte las presiones del mercado, un régimen intermedio \parencite{Boyer1978Optimal}.
\item \emph{Con agentes que optimizan}: en el modelo neokeynesiano de economía pequeña, lo óptimo es estabilizar la inflación interna y dejar flotar la moneda, y el tipo fijo es la peor de las reglas \parencite{GaliMonacelli2005Monetary} \vertema{3.A.32}.
\item \emph{Pero el miedo a flotar tiene razones}: si los precios de importación se fijan en la moneda del comprador o en dólares, el tipo de cambio no reasigna el gasto a corto plazo, y la política óptima debe atender también a sus desalineamientos \parencite{Engel2011Currency}.
\end{itemize}

\paragraph{d) Valoración}

\begin{itemize}
\item Explica bien por qué no hay un régimen óptimo para todos: depende de la estructura de las perturbaciones de cada país, que además cambia con el tiempo.
\item Pero es un modelo de precios fijos y de un solo objetivo: deja fuera la credibilidad, la inflación y el sistema financiero, que son los que más pesan en la evidencia.
\end{itemize}

\subsubsection{Evidencia empírica}\label{sec:evidencia-regimenes}

\begin{itemize}
\item \emph{El problema de medir}: lo que un país dice (\emph{de jure}) y lo que hace (\emph{de facto}) difieren, y a menudo el régimen de hecho es menos flexible que el declarado.\footnote{Por eso el FMI adoptó en 1999 una clasificación basada en lo que los países hacen: su informe anual recoge a la vez el régimen declarado y la evaluación de su personal \parencite{FMI2004Clasificacion}.} Por eso hay clasificaciones alternativas a la del FMI:
  \begin{itemize}
  \item \textcite{ReinhartRogoff2004Modern} usan el tipo de cambio del mercado paralelo y la inflación, y añaden la categoría de \emph{caída libre} (inflación de más del 40~\% anual); la ponen al día \textcite{IlzetzkiReinhartRogoff2019}, que muestran que muchos países que dicen flotar mantienen un ancla, sobre todo el dólar;
  \item \textcite{LevyYeyatiSturzenegger2005Classifying} agrupan los países por la volatilidad del tipo de cambio, de sus variaciones y de las reservas;
  \item \textcite{Shambaugh2004Effect} los clasifica por los movimientos frente a la moneda ancla.
  \end{itemize}
\item Según la clasificación, un mismo país es fijo o flotante, y por eso los estudios se contradicen: unos encuentran mejores resultados con los tipos fijos duros, otros con la flotación y otros con una flexibilidad limitada.
\item \emph{Resultados}:
  \begin{itemize}
  \item \emph{inflación}: los tipos fijos se asocian a menos inflación, por la credibilidad \parencite{GhoshGuldeWolf2002Exchange};
  \item \emph{crecimiento}: en los países en desarrollo, los regímenes más rígidos crecen menos y con más volatilidad; en los industriales no hay diferencia \parencite{LevyYeyatiSturzenegger2003Float};
  \item \emph{grado de desarrollo}: en los países con poco acceso a los mercados, un ancla da credibilidad sin coste de crecimiento; a medida que se desarrollan y se integran, gana la flexibilidad, y en los avanzados la flotación se asocia a más crecimiento sin perder credibilidad \parencite{RogoffEtAl2004Evolution};
  \item \emph{crisis}: los intermedios son los más propensos; los duros, a los colapsos del crecimiento \parencite{GhoshOstryQureshi2015}.
  \end{itemize}
\item \emph{Balance}: países parecidos eligen regímenes muy distintos sin grandes diferencias de crecimiento o de inflación, y sabemos poco de las causas y de las consecuencias de la elección \parencite{Rose2011Flaky}. La conclusión de \textcite{Frankel1999NoSingle} sigue en pie, y lo dice su título: ningún régimen es el adecuado para todos los países ni en todo momento.
\end{itemize}

\begin{anotaciones}
\begin{itemize}
\item Un repaso de la literatura sobre la elección del régimen: \parencite{Ouchen2013Optimal}.
\item La clasificación de hecho de \autor{Ethan Ilzetzki}, \autor{Carmen Reinhart} y \autor{Kenneth Rogoff}, con datos por país desde 1946: \parencite{IRR2026Datos}.
\end{itemize}
\end{anotaciones}

\begin{itemize}
\item Recapitulando el apartado: con libertad de capitales, cada régimen deja viva una sola política de demanda (el trilema), y la flotación no aísla del todo del ciclo financiero global; el tipo fijo da credibilidad y estabilidad al comercio a cambio de ajustes internos dolorosos y riesgo de crisis; y la elección depende de las perturbaciones, la integración con el ancla, la credibilidad y el desarrollo financiero de cada país.
\item Pero lo hemos analizado con 2 regímenes estilizados. En la práctica hay un menú de regímenes según el grado de fijación, la frecuencia con que se revisa el ancla y el compromiso que se anuncia. Pasamos a estudiarlo.
\end{itemize}

% ==========================================================================
\section{Implementación del régimen cambiario}\label{sec:implementacion}
\minutos{6}

\begin{itemize}
\item Una vez comparados los 2 regímenes estilizados, pasamos a los que los países tienen de verdad: el menú de regímenes que clasifica el FMI y hacia dónde se mueven los países.
\item \emph{Conclusión anticipada}: entre los 2 polos hay un continuo de regímenes, y la mayoría de los países está en medio. La visión bipolar no se ha cumplido: los regímenes intermedios sobreviven, pero solo los que se defienden con reservas, intervención o credibilidad.
\end{itemize}

% --------------------------------------------------------------------------
\subsection{Menú de regímenes cambiarios:\cortetitulo grado de fijación y publicidad}\label{sec:menu}

\begin{itemize}
\item El FMI publica desde 1950 un informe anual sobre los regímenes cambiarios y las restricciones de sus miembros, hoy el AREAER. Desde 1999 clasifica los regímenes por lo que los países hacen (\emph{de facto}), y desde 2009 con un sistema de 10 regímenes en 4 grupos \parencite{HabermeierEtAl2009Revised,FMI2024AREAER2023}.
\item La clasificación ordena los regímenes con 2 criterios:
  \begin{itemize}
  \item el \concepto{grado de fijación}: cuánto determina el tipo de cambio el mercado y cuánto la acción oficial;
  \item el \concepto{grado de publicidad}: el compromiso, formal o informal, de las autoridades con una senda del tipo de cambio.
  \end{itemize}
\item Cuanto más fijo y más público es el compromiso, más credibilidad da el régimen y menos autonomía monetaria deja; cuanto más flexible, al revés (gráfico~\ref{fig:regimenes-espectro}).
\end{itemize}

\begin{figure}[H]
\centering
\caption{El menú de regímenes cambiarios, del más rígido al más flexible}\label{fig:regimenes-espectro}
\grafico[0.92\textwidth]{regimenes-espectro}
\fuente{Elaboración propia a partir de \textcite{HabermeierEtAl2009Revised} y \textcite{FMI2024AREAER2023}. Porcentaje de los miembros del FMI en cada grupo, a 30 de abril de 2023.}
\end{figure}

\begin{itemize}
\item \emph{Cuántos hay en cada grupo}: según la clasificación de hecho del FMI, a 30 de abril de 2023, el 44,8~\% de sus miembros tenía un anclaje blando, el 32,5~\% flotaba, el 13,4~\% tenía un anclaje duro y el 9,3~\% otros regímenes gestionados \parencite{FMI2024AREAER2023} (tabla~\ref{tab:areaer}).
  \begin{itemize}
  \item La clasificación preliminar a 30 de abril de 2026 da un reparto parecido: el 47~\%, el 33~\%, el 14~\% y el 6~\% \parencite{FMI2026InformeAnual}.
  \end{itemize}
\item \emph{Las uniones monetarias no son un régimen aparte}: el FMI clasifica a los miembros de una unión por el régimen de la moneda común \vertema{3.B.16}.
  \begin{itemize}
  \item Los países del euro aparecen en libre flotación: el euro flota frente al resto del mundo, aunque entre ellos no haya tipo de cambio.
  \item Las 2 uniones del franco CFA, la UEMOA y la CEMAC, tienen un anclaje convencional al euro, y la del Caribe Oriental, una caja de conversión con el dólar \parencite{FMI2024AREAER2023}.
  \end{itemize}
\end{itemize}

\begin{table}[H]
\centering
\caption{Los regímenes cambiarios según el FMI}\label{tab:areaer}
\small
\begin{tabular}{>{\raggedright\arraybackslash}m{2.0cm}>{\raggedright\arraybackslash}m{2.6cm}>{\raggedright\arraybackslash}m{5.0cm}>{\centering\arraybackslash}m{1.2cm}>{\raggedright\arraybackslash}m{3.4cm}}
\toprule
\rowcolor{tceeazulpalido}\textbf{Grupo} & \textbf{Régimen} & \textbf{Criterio} & \textbf{\% (2023)} & \textbf{Ejemplos (2026)}\\
\midrule
\multirow{2}{=}{\textcolor{tceeazul}{\textbf{Anclaje duro}} (13,4~\%)} & Sin moneda propia de curso legal & Usa como única moneda de curso legal la de otro país & 7,2 & Ecuador, Panamá (dólar); Montenegro, San Marino (euro)\\
 & Caja de conversión & Paridad fijada por ley y base monetaria respaldada con divisas & 6,2 & Hong Kong, Bosnia y Herzegovina, Caribe Oriental\\
\midrule
\multirow{5}{=}{\textcolor{tceedorado}{\textbf{Anclaje blando}} (44,8~\%)} & Convencional & Fluctúa menos del $\pm1$~\% en torno a una paridad (o en un margen del 2~\%) durante 6 meses & 20,6 & Arabia Saudí, Dinamarca, UEMOA, CEMAC\\
 & Estabilizado & Margen del 2~\% durante 6 meses por acción oficial, sin compromiso & 11,3 & Macedonia del Norte, Bolivia, Nicaragua\\
 & Deslizamiento & Paridad ajustada en pequeñas cantidades, a un ritmo fijo o según indicadores & 1,5 & Botsuana\\
 & Similar al deslizamiento & Margen del 2~\% en torno a una tendencia durante 6 meses & 10,8 & China, Singapur, Vietnam\\
 & Bandas horizontales & Banda de al menos $\pm1$~\% en torno a una paridad & 0,5 & Marruecos\\
\midrule
\multirow{2}{=}{\textcolor{tceerojo}{\textbf{Flotación}} (32,5~\%)} & Flotación & Lo fija sobre todo el mercado; se interviene para suavizar & 16,5 & Brasil, Nueva Zelanda, Corea, Egipto\\
 & Flotación libre & Intervención excepcional: como mucho 3 veces en 6 meses, de no más de 3 días hábiles & 16,0 & Estados Unidos, zona del euro, Japón, Reino Unido\\
\midrule
\textcolor{tceegris}{\textbf{Residual}} (9,3~\%) & Otros gestionados & No encaja en los anteriores & 9,3 & Argentina, Irán, Venezuela\\
\bottomrule
\end{tabular}
\fuente{\textcite{FMI2024AREAER2023} (porcentajes de los 190 miembros y 4 territorios, a 30 de abril de 2023), \textcite{FMI2026InformeAnual} (ejemplos, clasificación preliminar a 30 de abril de 2026) y \textcite{HabermeierEtAl2009Revised} (criterios). Los países del euro cuentan en libre flotación; los de la UEMOA y la CEMAC, en el anclaje convencional.}
\end{table}

\begin{notaopositor}
En el cante basta con los 4 grupos, su porcentaje y un ejemplo de cada uno; los 10 regímenes, uno a uno y con sus series, están en el anexo~\ref{anexo:regimenes}.
\end{notaopositor}

\subsubsection{Los anclajes duros}\label{sec:anclajes-duros}

\begin{itemize}
\item En un \concepto{anclaje duro} (\emph{hard peg}), la autoridad fija la paridad con un compromiso institucional que hace muy caro abandonarla. A cambio, renuncia por completo a la política monetaria propia.
\item \emph{Sin moneda propia de curso legal}: el país usa como única moneda de curso legal la de otro país, la \concepto{dolarización oficial}.
  \begin{itemize}
  \item Puede ser con acuerdo con el emisor (San Marino y Andorra con la UE) o unilateral (Kosovo y Montenegro, que usan el euro sin acuerdo; Ecuador, El Salvador y Panamá, el dólar).\footnote{No hay que confundirla con la \emph{dolarización no oficial} (o \emph{de facto}): el uso de una divisa en los depósitos, los préstamos o los precios sin que sea de curso legal. El Salvador dio curso legal al bitcóin junto al dólar en septiembre de 2021; con la reforma de enero de 2025, pactada con el FMI, aceptarlo pasó a ser voluntario \parencite{FMI2024ElSalvadorSLA,FMI2025ElSalvador}. Para el episodio, \parencite{Funcas2022ODF77}.}
  \item \emph{Ventajas}: más credibilidad que cualquier tipo fijo, porque no hay moneda que devaluar, y más estabilidad de precios.
  \item \emph{Inconvenientes}: el país importa la política monetaria de un banco central que no mira su ciclo, pierde el señoreaje y no tiene prestamista de última instancia; a diferencia de una unión monetaria, no influye en la política común.
  \end{itemize}
\item \emph{La unión monetaria}: los países comparten moneda y banco central, y cada uno pesa en la política común. Es el extremo de los anclajes duros \vertema{3.B.16}.
\item \emph{La caja de conversión} (\emph{currency board}): fija la paridad por ley y obliga a respaldar con divisas toda la base monetaria.
  \begin{itemize}
  \item En el balance del banco central, $\mathrm{BM}=\mathrm{CI}+R$: la base monetaria tiene como contrapartida el crédito interno y las reservas. En una caja de conversión el crédito interno no crece y $R\geq\mathrm{BM}$, así que la base solo aumenta si entran divisas.
  \item La autoridad no hace operaciones de mercado abierto, no presta al Gobierno ni garantiza a los bancos. Solo puede prestar el exceso de reservas sobre la base: no es un prestamista de última instancia.
  \item \emph{Ventajas}: una credibilidad casi de dolarización, conservando la moneda propia y algo de señoreaje (los intereses de las reservas).
  \item \emph{Inconvenientes}: sin política monetaria y con el ajuste por precios y salarios. Si la inflación nace de las expectativas, la caja la frena; si nace de causas estructurales (por ejemplo, sindicatos con mucho poder de negociación), apenas sirve.
  \item \emph{Ejemplos}: Hong Kong mantiene desde 1983 una caja de conversión con el dólar estadounidense, y desde 2005 su autoridad monetaria se compromete a comprar y vender dólares a 7,75 y 7,85 dólares de Hong Kong \parencite{HKMA2026LERS,HKMA2005Refinamientos}; Estonia, Lituania y Bulgaria la tuvieron hasta entrar en el euro, Bulgaria el 1 de enero de 2026 \vertema{3.B.16}.
  \item \emph{El límite}: Argentina (1991-2001) fijó el peso al dólar con una caja de conversión; con el peso sobrevalorado, deuda en dólares y una política fiscal que no la sostenía, acabó en el corralito y la devaluación de 2002 \vertema{3.B.17} \vertema{3.B.21}.
  \end{itemize}
\end{itemize}

\begin{anotaciones}
\textbf{Cuándo funciona una caja de conversión}
\begin{itemize}
\item La moneda ancla es estable y su país está integrado con el propio: ciclos parecidos y mucho comercio entre ellos.
\item La paridad es adecuada, y la autoridad tiene un margen de reservas por encima del 100~\% de la base para atender contingencias sin poner en riesgo la convertibilidad.
\item La caja es independiente del Gobierno, y el sistema bancario es sólido, porque no tendrá prestamista de última instancia.
\end{itemize}
\end{anotaciones}

\subsubsection{Los anclajes blandos}\label{sec:anclajes-blandos}

\begin{itemize}
\item En un \concepto{anclaje blando} (\emph{soft peg}), la autoridad defiende una paridad, una senda o una banda con intervención y con los tipos de interés, pero sin un compromiso institucional que impida cambiarla. Da algo de credibilidad y deja algo de margen, y por eso es vulnerable a los ataques.
\item \emph{Convencional}: la moneda se ancla a otra o a una cesta y fluctúa menos del $\pm1$~\%.
  \begin{itemize}
  \item Arabia Saudí mantiene el riyal en 3,75 por dólar desde 1986 \parencite{SAMA2016TipoCambio,BIS2026TiposBilaterales}; Dinamarca ancla la corona al euro dentro del mecanismo de tipos de cambio europeo (MTC~II) \vertema{3.B.43}.
  \end{itemize}
\item \emph{Estabilizado}: el tipo se mantiene en un margen del 2~\% durante 6 meses por acción oficial, sin compromiso anunciado.
  \begin{itemize}
  \item Macedonia del Norte declara un régimen de flotación, pero mantiene el denar en torno a 61,5 por euro desde hace más de 20 años \parencite{FMI2026InformeAnual,BIS2026TiposBilaterales}.
  \end{itemize}
\item \emph{Deslizamiento} (\emph{crawling peg}): la paridad se ajusta en pequeñas cantidades, a un ritmo anunciado o según indicadores como la inflación pasada.
  \begin{itemize}
  \item Sirve a un país con más inflación que el ancla: el tipo nominal se deprecia poco a poco y el real no se aprecia. Pero deja poco margen monetario, y no elimina el riesgo de apreciación real ni de crisis, como mostró México en 1994-1995 \vertema{3.B.17}.
  \item Nicaragua depreció el córdoba un poco cada día durante 30 años, a un ritmo anunciado (el 5~\% anual hasta 2019, y menos después), hasta fijarlo en 36,6243 por dólar el 1 de enero de 2024 \parencite{BCN2023Deslizamiento,BIS2026TiposBilaterales}. Hoy el único deslizamiento que reconoce el FMI es el de Botsuana \parencite{FMI2026InformeAnual}.
  \end{itemize}
\item \emph{Similar al deslizamiento} (\emph{crawl-like}): el tipo se mueve en un margen del 2~\% alrededor de una tendencia, sin compromiso anunciado.
  \begin{itemize}
  \item China, con un régimen oficial de flotación gestionada frente a una cesta, figura de hecho en este grupo \parencite{FMI2026InformeAnual}.
  \end{itemize}
\item \emph{Bandas horizontales}: el tipo puede moverse en una banda de al menos el $\pm1$~\% en torno a una paridad central. Hoy solo la tiene Marruecos, pero fue el régimen de los puntos del oro, de Bretton Woods (con bandas del $\pm1$~\%) y del Sistema Monetario Europeo.
  \begin{itemize}
  \item Una banda es tanto más parecida a un tipo fijo cuanto más estrecha, más pública y más firme sea, y cuanto menos se realinee: el SME de los primeros 80, con realineamientos frecuentes, funcionó casi como un deslizamiento.
  \item Ensancharla tiene 2 efectos opuestos: deja más volatilidad, pero también más margen a la política monetaria y menos ocasiones de tocar los bordes, lo que puede dar más credibilidad, como en el SME tras ensanchar las bandas en 1993.
  \item La credibilidad es la clave, y la estudia el modelo de las zonas objetivo de \textcite{Krugman1991Target}.
  \end{itemize}
\end{itemize}

\subsubsection{Las bandas:\cortetitulo zonas objetivo de \autor{Krugman}}\label{sec:krugman}

\paragraph{a) Idea}

\begin{itemize}
\item ¿Estabiliza una banda creíble el tipo de cambio más de lo que lo haría la intervención por sí sola? \textcite{Krugman1991Target} muestra que sí: las expectativas trabajan a favor del banco central.
\item El modelo dio a la credibilidad, que no es una variable macroeconómica, un papel central en el análisis de los regímenes.
\end{itemize}

\paragraph{b) Modelo}

\subparagraph{i) Supuestos}

\begin{itemize}
\item El enfoque monetario con precios flexibles y paridad del poder adquisitivo \vertema{3.B.14}, en tiempo continuo y en logaritmos, con $e_t=\ln E_t$: el tipo de cambio depende de los fundamentos y de la depreciación esperada,
  \begin{equation}\label{eq:zona-objetivo}
  e_t=m_t+v_t+\alpha\,\frac{\mathbb{E}_t[\mathrm{d}e_t]}{\mathrm{d}t},
  \end{equation}
  con $m_t$ la oferta monetaria, $v_t$ una perturbación de la demanda de dinero (la renta y la velocidad) que sigue un paseo aleatorio y $\alpha>0$ la semielasticidad de la demanda de dinero respecto al tipo de interés.
\item La autoridad fija una banda, $e^{\min}\leq e_t\leq e^{\max}$, y solo interviene en los bordes: si la moneda toca el borde depreciado, reduce $m_t$; si toca el apreciado, lo aumenta; dentro de la banda, $m_t$ no cambia.
\item La banda es creíble: el mercado sabe que la autoridad la defenderá.
\end{itemize}

\subparagraph{ii) Desarrollo}

\begin{itemize}
\item \emph{Sin banda}, el tipo de cambio seguiría a los fundamentos: es la recta de 45° del gráfico~\ref{fig:zona-objetivo}, el tipo de cambio sombra, $\tilde e=m+v$.
\item \emph{Con una banda creíble}, el tipo no se mueve libremente dentro de ella.
  \begin{itemize}
  \item Cerca del borde depreciado, el mercado sabe que la autoridad impedirá que la moneda se deprecie más: espera una apreciación, y compra la moneda, que se aprecia ya.
  \item Cerca del borde apreciado ocurre lo contrario.
  \item La relación entre los fundamentos y el tipo de cambio es una \emph{S} aplanada, menos inclinada que la de 45° y tangente a los bordes (\emph{smooth pasting}) (anexo~\ref{anexo:krugman}).
  \end{itemize}
\end{itemize}

\begin{figure}[H]
\centering
\caption{Las zonas objetivo de \autor{Krugman}: luna de miel y divorcio}\label{fig:zona-objetivo}
\grafico[0.66\textwidth]{zona-objetivo}
\fuente{Elaboración propia a partir de \textcite{Krugman1991Target} y \textcite{BertolaCaballero1992Target}. La recta de 45° es el tipo de cambio sombra; la \emph{S} aplanada, la banda creíble; la \emph{S} invertida, la banda sin credibilidad; $\pm\bar f$, los bordes de los fundamentos, donde la autoridad interviene.}
\end{figure}

\subparagraph{iii) Implicaciones}

\begin{itemize}
\item Es el \concepto{efecto luna de miel}: una banda creíble estabiliza por sí sola, el tipo de cambio pasa mucho tiempo lejos de los bordes y la intervención necesaria es mínima.
\item Lo que mueve hoy el tipo de cambio es lo que se espera de él mañana.
\end{itemize}

\paragraph{c) Extensiones}

\begin{itemize}
\item \emph{Una banda sin credibilidad}: si el mercado cree que, al llegar al borde, la autoridad realineará en vez de defenderlo, la expectativa se da la vuelta \parencite{BertolaCaballero1992Target}.
  \begin{itemize}
  \item Cerca del borde depreciado se espera más depreciación, y los agentes venden la moneda: la depreciación se adelanta.
  \item La curva se vuelve una \emph{S} invertida, más inclinada que la de 45°, y solo la paridad central es estable.
  \item La banda añade volatilidad en vez de quitarla: es lo que suele llamarse \concepto{efecto divorcio}.
  \end{itemize}
\end{itemize}

\paragraph{d) Evidencia empírica}

\begin{itemize}
\item Los datos del SME no se parecían a la \emph{S} creíble: el tipo de cambio pasaba mucho tiempo cerca de los bordes, y los diferenciales de tipos de interés anticipaban los realineamientos \parencite{BertolaCaballero1992Target}.
\item Volvamos a España: la peseta entró en el mecanismo de cambios del SME en 1989; tras la salida de la libra en septiembre de 1992, se devaluó varias veces entre 1992 y 1993, y en agosto de 1993 las bandas se ensancharon al $\pm15$~\% \parencite{EichengreenWyplosz1993Unstable} \vertema{3.B.17} \vertema{3.B.43}. Fue el divorcio de una banda que el mercado dejó de creer.
\end{itemize}

\paragraph{e) Valoración}

\begin{itemize}
\item Lo que estabiliza no es la banda, sino la credibilidad. Una banda creíble funciona sola; una que no lo es invita al ataque.
\item El modelo supone precios flexibles y fundamentos exógenos; no explica por qué una autoridad abandona la banda, que es el terreno de las crisis de segunda generación \vertema{3.B.17}.
\end{itemize}

\subsubsection{La flotación y los otros regímenes gestionados}\label{sec:flotacion}

\begin{itemize}
\item \emph{Flotación}: el mercado fija el tipo de cambio, pero la autoridad interviene para suavizar sus movimientos sin oponerse a las tendencias de fondo, «ir contra el viento».
  \begin{itemize}
  \item \emph{Ventajas}: deja autonomía monetaria y no ofrece a los especuladores una apuesta en un solo sentido.
  \item \emph{Inconvenientes}: puede no anclar las expectativas de inflación, porque la depreciación es posible, y obliga a cubrir el riesgo de cambio.
  \item \emph{Ejemplos}: Brasil y Nueva Zelanda; Egipto, desde que en marzo de 2024 unificó sus tipos y dejó flotar la libra \parencite{FMI2026InformeAnual,CBE2024Flotacion}.
  \end{itemize}
\item \emph{Flotación libre}: la intervención es excepcional, para corregir desórdenes del mercado, como mucho 3 veces en 6 meses y de no más de 3 días hábiles cada vez \parencite{HabermeierEtAl2009Revised}. Así flotan Estados Unidos, la zona del euro, Japón o el Reino Unido.
\item \emph{Otros regímenes gestionados}: la categoría residual, que recoge lo que no encaja en las demás, a menudo regímenes que cambian con frecuencia. Ahí figuran hoy Argentina, Irán y Venezuela \parencite{FMI2026InformeAnual}.
\item Recapitulando el menú, cada régimen compra credibilidad con autonomía, y los intermedios son los más vulnerables a los ataques (tabla~\ref{tab:intermedios}).
\item Volvamos a España, que ha recorrido casi todo el menú: tipos de cambio múltiples hasta 1959 (\ref{sec:regulacion}), la paridad ajustable de Bretton Woods, las bandas del SME desde 1989 y, desde 1999, el euro: una unión monetaria con sus socios que flota libremente frente al resto del mundo \vertema{3.B.16} \vertema{4.A.27}.
\end{itemize}

\begin{table}[H]
\centering
\caption{Los regímenes, del más rígido al más flexible}\label{tab:intermedios}
\small
\begin{tabular}{>{\raggedright\arraybackslash}m{3.0cm}>{\centering\arraybackslash}m{2.3cm}>{\centering\arraybackslash}m{2.2cm}>{\centering\arraybackslash}m{2.3cm}>{\raggedright\arraybackslash}m{4.2cm}}
\toprule
\rowcolor{tceeazulpalido}\textbf{Régimen} & \textbf{Credibilidad} & \textbf{Autonomía monetaria} & \textbf{Vulnerable a ataques} & \textbf{Para quién}\\
\midrule
Unión monetaria, dolarización & \textcolor{tceeverde}{Máxima} & \textcolor{tceerojo}{Nula} & \textcolor{tceeverde}{No} (salir es carísimo) & Muy integrado con el ancla, o con un historial de inflación\\
Caja de conversión & \textcolor{tceeverde}{Muy alta} & \textcolor{tceerojo}{Nula} & \textcolor{tceedorado}{Poco} & Pequeño, abierto y con una banca sólida\\
Convencional y bandas estrechas & \textcolor{tceedorado}{Media} & \textcolor{tceerojo}{Escasa} & \textcolor{tceerojo}{Mucho} & Con controles o muchas reservas\\
Deslizamiento y bandas móviles & \textcolor{tceedorado}{Media} & \textcolor{tceedorado}{Algo} & \textcolor{tceerojo}{Bastante} & Con más inflación que el ancla\\
Flotación & \textcolor{tceedorado}{La del banco central} & \textcolor{tceeverde}{Alta} & \textcolor{tceeverde}{Poco} & Emergentes con objetivo de inflación\\
Flotación libre & \textcolor{tceedorado}{La del banco central} & \textcolor{tceeverde}{Plena} & \textcolor{tceeverde}{No} & Grandes, con mercados profundos\\
\bottomrule
\end{tabular}
\fuente{Elaboración propia a partir de \textcite{Frankel1999NoSingle}, \textcite{Fischer2001Bipolar} y \textcite{GhoshOstryQureshi2015}. En verde, ventaja; en dorado, intermedio; en rojo, inconveniente.}
\end{table}

% --------------------------------------------------------------------------
\subsection{¿Hacia qué régimen cambiario tienden los países?}\label{sec:tendencias}

\begin{itemize}
\item Conocido el menú, la pregunta es qué eligen los países en la práctica. En la segunda mitad de los 90 había 2 predicciones enfrentadas.
\item \emph{La visión bipolar}: con libertad de capitales, los regímenes intermedios desaparecerían y solo quedarían los extremos, la flotación o el anclaje duro \parencite{Eichengreen1994International,ObstfeldRogoff1995Mirage}. Tras las crisis de los 90, \textcite{Fischer2001Bipolar} vio el «vaciado» de los intermedios. Pegatina para los intermedios: «ni chicha ni limonada».
\item \emph{El miedo a flotar}: muchos países que dicen flotar no lo hacen; sus tipos de cambio se mueven poco y sus reservas y tipos de interés, mucho \parencite{CalvoReinhart2002}.
  \begin{itemize}
  \item ¿Por qué seguir una banda en secreto, si su virtud es que, siendo creíble, estabiliza sola? Hay 2 razones enfrentadas.
  \item \emph{A favor de anunciarla}: si la moneda va a estar en una banda de todos modos, es mejor que sea pública y fija, para ganar credibilidad contra la inflación y la volatilidad.
  \item \emph{A favor de callarla}: una banda que no se anuncia no ofrece un blanco a los ataques especulativos ni obliga a un ajuste interno costoso para defenderla.
  \end{itemize}
\item \emph{Lo que ha pasado}: ni lo uno ni lo otro del todo (gráfico~\ref{fig:regimenes-evolucion}).
  \begin{itemize}
  \item Los anclajes blandos siguen siendo el grupo más numeroso, con cerca de la mitad de los países, y el reparto entre los 4 grupos apenas ha cambiado en 10 años \parencite{FMI2024AREAER2023,FMI2026InformeAnual}.
  \item Con su clasificación propia, \textcite{IlzetzkiReinhartRogoff2019} encuentran que los regímenes de flexibilidad limitada son mayoría y que el dólar es, con mucho, el ancla dominante.
  \end{itemize}
\end{itemize}

\begin{figure}[H]
\centering
\caption{Los regímenes cambiarios de los miembros del FMI, 2015-2026 (\%)}\label{fig:regimenes-evolucion}
\grafico[0.8\textwidth]{regimenes-evolucion}
\fuente{Elaboración propia con datos del FMI \parencite{FMI2024AREAER2023,FMI2026InformeAnual}. Porcentaje de los miembros a 30 de abril de cada año; 2026, clasificación preliminar; 2024 y 2025, sin dato publicado.}
\end{figure}

\begin{itemize}
\item \emph{Los hechos estilizados}, según \textcite{Rose2011Flaky} y \textcite{KleinShambaugh2010Exchange}:
  \begin{itemize}
  \item solo un tercio de los países flota (65 de 195 economías en 2026), y casi todos los que flotan tienen un objetivo de inflación \parencite{FMI2026InformeAnual};
  \item muchos más fijan, casi siempre al dólar (sobre todo los exportadores de petróleo y los centros financieros) y algunos al euro;
  \item los que fijan son, salvo China, pequeños: pesan poco en el PIB mundial. Los países ricos y grandes flotan, pero no hay una relación clara entre la renta per cápita y el régimen;
  \item los regímenes duran cada vez más.
  \end{itemize}
\item \emph{Una perspectiva histórica}:
  \begin{itemize}
  \item en los 80 dominaba la idea de fijar: los planes de estabilización de América Latina tras la crisis de la deuda y, en Europa, el SME camino de la unión monetaria;
  \item entre 1992 (la crisis del SME) y 2002 (el fin de la caja de conversión argentina), las crisis cambiarias graves llevaron a muchos países a no anunciar tipos fijos y a seguir bandas en la práctica: el miedo a flotar \vertema{3.B.17};
  \item desde entonces no ha habido crisis cambiarias de esa magnitud, ni siquiera con la Gran Recesión, en parte porque los emergentes flotan con intervención y han acumulado reservas (\ref{sec:intervencion}).
  \end{itemize}
\item \emph{Las unificaciones recientes}: unificar un mercado cambiario dividido suele exigir una devaluación grande.
  \begin{itemize}
  \item Egipto dejó flotar la libra en marzo de 2024 para cerrar la brecha con el mercado paralelo \parencite{CBE2024Flotacion}.
  \item Argentina levantó en abril de 2025 casi todo el cepo cambiario y pasó a flotar entre bandas \parencite{BCRA2025Fase3,FMI2025ArgentinaEFF}.
  \end{itemize}
\item Recapitulando el apartado: el menú va de la unión monetaria a la flotación libre; la mitad de los países vive en medio, y los intermedios que sobreviven son los que tienen credibilidad o reservas para defenderse.
\item Para sostener cualquier régimen que no sea la flotación libre hay que intervenir en el mercado de divisas o regular los movimientos de capitales. Son las 2 herramientas que estudiamos a continuación.
\end{itemize}

% ==========================================================================
\section{Intervención y regulación de los mercados de cambio}\label{sec:intervencion-regulacion}
\minutos{10}

\begin{itemize}
\item Elegido el régimen, hay que sostenerlo. Pasamos a las 2 herramientas con que las autoridades actúan sobre el mercado de divisas: la intervención, que cambia la oferta y la demanda de divisas, y la regulación, que limita quién puede comprarlas y venderlas y para qué.
\item \emph{Conclusión anticipada}: las 2 funcionan, pero con límites. La intervención suaviza los movimientos y funciona mejor si es creíble y pública, pero no sostiene un tipo de cambio contra los fundamentos. Los controles de capitales cuestan bienestar en un mundo sin fricciones, pero pueden mejorarlo si corrigen una externalidad; cambian la composición de los flujos más que su volumen.
\end{itemize}

% --------------------------------------------------------------------------
\subsection{Intervención}\label{sec:intervencion}

\subsubsection{Definición, motivos y formas}\label{sec:intervencion-formas}

\begin{itemize}
\item \concepto{Intervenir} es comprar o vender divisas a cambio de moneda nacional (o anunciar que se hará) para influir en el tipo de cambio: las autoridades alteran el tipo que saldría de una flotación plenamente libre \vertema{3.B.13}.
\item \emph{Los motivos} dependen del régimen \vertema{3.B.11} \vertema{3.B.13}:
  \begin{itemize}
  \item con tipo fijo o bandas, la autoridad interviene para \emph{defender la paridad};
  \item con tipo flexible, para \emph{suavizar} los movimientos bruscos y corregir desórdenes del mercado («ir contra el viento»);
  \item también por \emph{precaución}: comprar divisas para acumular reservas, un seguro frente a las paradas súbitas de los flujos de capital que da confianza a los inversores;
  \item o para \emph{frenar una apreciación} y ganar competitividad. El FMI lo considera manipulación si busca una ventaja injusta (artículo~IV del Convenio Constitutivo) \vertema{3.B.19};
  \item los exportadores de materias primas acumulan cuando los precios están altos para gastar cuando bajan, a veces a través de fondos soberanos.
  \end{itemize}
\end{itemize}

\begin{anotaciones}
\textbf{Para qué sirven las reservas}
\begin{itemize}
\item \emph{Intervenir y dar confianza}: influir en el tipo de cambio y convencer a los inversores de que el país podrá pagar en divisas, lo que abarata su financiación y hace menos probable una crisis.
\item \emph{Pagar al exterior aunque se corte la financiación}: las importaciones, la deuda pública en divisas y la liquidez de emergencia en divisas a los bancos.
\item \emph{Gestionar la liquidez}: como garantía en los \emph{swaps} de divisas de la política monetaria y para que los mercados de divisas sean líquidos.
\item \emph{Invertir el exceso}: lo que sobra de lo necesario se invierte para obtener rentabilidad, a veces en fondos soberanos que financian programas de desarrollo.
\item Cuántas reservas bastan se mide con varios indicadores y con la métrica del FMI \vertema{3.B.22}.
\end{itemize}
\end{anotaciones}

\begin{itemize}
\item \emph{Las formas}: la intervención clásica es la compraventa de divisas, que se ve en el balance simplificado del banco central, $\mathrm{BM}=\mathrm{CI}+R$ (tabla~\ref{tab:balance-intervencion}).
  \begin{itemize}
  \item En la \concepto{intervención no esterilizada}, comprar divisas sube $R$ y la base monetaria en la misma cuantía: es una expansión monetaria pagada en divisas, y funciona por el canal de siempre, los tipos de interés. Con tipo fijo y capital libre, es la única que sostiene la paridad.
  \item En la \concepto{intervención esterilizada}, el banco central \emph{vende} a la vez títulos en moneda nacional (o emite los suyos) por el mismo importe: $\Delta\mathrm{CI}=-\Delta R$ y la base no cambia. Solo cambia la composición de las carteras privadas, con más bonos nacionales y menos activos en divisas \vertema{3.B.11}.
  \end{itemize}
\end{itemize}

\begin{table}[H]
\centering
\caption{Una compra de 100 de divisas en el balance del banco central}\label{tab:balance-intervencion}
\small
\begin{tabular}{>{\raggedright\arraybackslash}m{3.4cm}>{\centering\arraybackslash}m{4.4cm}>{\centering\arraybackslash}m{3.4cm}}
\toprule
\rowcolor{tceeazulpalido}\textbf{Intervención} & \textbf{Activo} & \textbf{Pasivo}\\
\midrule
No esterilizada & $\Delta R=+100$ & \textcolor{tceerojo}{$\Delta\mathrm{BM}=+100$}\\
Esterilizada & $\Delta R=+100$; $\Delta\mathrm{CI}=-100$ (venta de títulos) & \textcolor{tceeverde}{$\Delta\mathrm{BM}=0$}\\
\bottomrule
\end{tabular}
\fuente{Elaboración propia. $\mathrm{BM}$, base monetaria; $\mathrm{CI}$, crédito interno; $R$, reservas internacionales.}
\end{table}

\begin{itemize}
\item \emph{Otras formas}:
  \begin{itemize}
  \item la \concepto{persuasión} (\emph{jawboning}): la autoridad anuncia que intervendrá en una dirección, y el anuncio basta para mover las expectativas. Pegatina: «más vale una palabra a tiempo que cien a destiempo»;
  \item la \emph{intervención coordinada} de varios bancos centrales, como los acuerdos del Plaza (1985), para depreciar el dólar, y del Louvre (1987), para frenar su caída \parencite{DominguezFrankel1993Does} \vertema{3.B.19} \vertema{3.B.21};
  \item hoy también los derivados y las operaciones a plazo, y las líneas \emph{swap} entre bancos centrales, que dan divisas sin gastar reservas \vertema{3.B.43}.
  \end{itemize}
\item \emph{Un ejemplo reciente}: Japón compró yenes por unos 9,8 billones de yenes a finales de abril y principios de mayo de 2024, y por otros 5,5 billones en julio, para frenar la caída de su moneda a mínimos de 34 años frente al dólar \parencite{MOF2024Intervencion}.
\end{itemize}

\subsubsection{Canales y eficacia de la intervención esterilizada}\label{sec:intervencion-canales}

\begin{itemize}
\item La intervención no esterilizada funciona como cualquier política monetaria. La pregunta es si funciona la esterilizada, que no cambia la cantidad de dinero.
\item Con activos nacionales y extranjeros sustitutos perfectos, no: la prima de riesgo es nula en la paridad de intereses, $\rho=0$ en~\eqref{eq:trilema}, y la intervención no cambia ni los tipos de interés ni las expectativas.
\item Puede funcionar por 3 canales:
  \begin{itemize}
  \item el \concepto{canal de cartera}: si los activos no son sustitutos perfectos, la prima $\rho$ crece con la oferta de bonos nacionales en manos privadas. Al venderlos y comprar divisas, el banco central sube $\rho$ y deprecia la moneda sin tocar el dinero \parencite{Branson1977Asset} \vertema{3.B.14}. Con intermediarios de balance limitado, la intervención funciona porque alguien tiene que absorber los flujos \parencite{GabaixMaggiori2015};
  \item el \concepto{canal de señales}: la intervención anuncia la política monetaria futura o el tipo de cambio que el banco central cree de equilibrio, y el banco central «pone su dinero» detrás del anuncio; mueve las expectativas, $\gamma E^e$, aunque los activos sean sustitutos perfectos \parencite{Mussa1981Role};
  \item el \concepto{canal de coordinación}: cuando el tipo de cambio se aleja mucho de los fundamentos, la intervención ayuda a que los operadores que lo saben actúen a la vez \parencite{SarnoTaylor2001Official}; y como en el mercado el tipo se mueve con el flujo de órdenes, la del banco central lo mueve como cualquier otra \parencite{EvansLyons2002} \vertema{3.B.14}.
  \end{itemize}
\item \emph{Evidencia empírica}:
  \begin{itemize}
  \item en los 80 se pensaba que la esterilizada era inútil; las intervenciones coordinadas del Plaza y del Louvre, que movieron el dólar, cambiaron la opinión: la intervención funciona, sobre todo si es pública y coordinada \parencite{DominguezFrankel1993Does};
  \item con datos diarios de 33 países, la intervención tiene éxito según algunos criterios en más del 80~\% de los casos, sobre todo para suavizar y con bandas estrechas, y más si se anuncia y se acompaña de comunicación \parencite{FratzscherEtAl2019When};
  \item en los emergentes, la intervención frena la apreciación ante las entradas de capital \parencite{BlanchardAdlerCarvalho2015Can}.
  \end{itemize}
\item \emph{Valoración}: la intervención mueve el tipo de cambio a corto plazo y funciona mejor para frenar movimientos bruscos que para cambiar tendencias. No puede sostener un tipo de cambio contra los fundamentos: vender reservas para defender una paridad sobrevalorada acaba en crisis \vertema{3.B.17}.
\end{itemize}

\subsubsection{Costes, reservas y el Marco de Política Integrado}\label{sec:intervencion-costes}

\begin{itemize}
\item \emph{El coste cuasifiscal de esterilizar}: el banco central paga por los títulos que vende más de lo que rinden las reservas que compra, y mantiene altos los tipos de interés que atraen más capital \parencite{Calvo1991Esterilizacion}.
\item \emph{El coste de oportunidad de las reservas}: son recursos que podrían invertirse en proyectos más rentables; para los países en desarrollo, en torno al 1~\% del PIB al año \parencite{Rodrik2006Social}.
\item \emph{Las reservas se han disparado}: las de divisas de los bancos centrales suman 13,2 billones de dólares (junio de 2026), frente a 1,8 billones en 1999; el 56,7~\% de las asignadas está en dólares y el 20,6~\% en euros \parencite{FMI2026COFER}, y China tiene unos 3,4 billones \parencite{SAFE2026Reservas} (gráfico~\ref{fig:reservas-mundiales}) \vertema{3.B.11}.
  \begin{itemize}
  \item Fueron la consecuencia de intervenir para frenar la apreciación, como China para mantener el renminbi casi fijo al dólar, y de la lección de la crisis asiática: ante un mercado que se mueve en rebaño, cualquier volumen de reservas parecía poco.
  \end{itemize}
\end{itemize}

\begin{figure}[H]
\centering
\caption{Las reservas de divisas de los bancos centrales, 1999-2026}\label{fig:reservas-mundiales}
\grafico[0.8\textwidth]{reservas-mundiales}
\fuente{Elaboración propia con datos del FMI \parencite{FMI2026COFER}. Total en billones de dólares y porcentaje de las reservas asignadas que está en dólares; datos a final de año y, en 2026, de junio.}
\end{figure}

\begin{itemize}
\item \emph{La política monetaria como intervención implícita}: las compras de activos de los bancos centrales de moneda fuerte deprecian su moneda sin tocar las reservas \parencite{DedolaEtAl2021Bazooka} \vertema{3.B.13}.
  \begin{itemize}
  \item No es intervención en sentido estricto, porque no se gestionan reservas, pero en espíritu cambia el valor de la moneda, y por eso tensa las relaciones entre países \vertema{3.B.19}.
  \item En enero de 2018, \autor{Mario Draghi} respondió al secretario del Tesoro estadounidense, \autor{Steven Mnuchin}, que había defendido en Davos un dólar débil \parencite{BCE2018Draghi,Bloomberg2018DraghiMnuchin}.
  \end{itemize}
\item \emph{El Marco de Política Integrado del FMI}:
  \begin{itemize}
  \item En \autor{Mundell}-\autor{Fleming}, la flotación y la política monetaria bastan para aislar a un país. Con fricciones (mercados de divisas poco profundos, descalces de moneda, comercio facturado en dólares, expectativas de inflación mal ancladas), no.
  \item El \concepto{Marco de Política Integrado} estudia juntos 4 instrumentos: la política monetaria, la intervención cambiaria, las medidas macroprudenciales y las de gestión de los flujos de capital \parencite{BasuEtAl2020Conceptual,FMI2020IPF}.
  \item Cada instrumento sirve para una fricción: un instrumento más no siempre conviene, porque puede no ser el adecuado para la imperfección que hay \parencite{BasuEtAl2020Conceptual}.
  \item Según los principios que aprobó el FMI en 2023, la intervención solo está justificada ante fricciones identificadas (mercados poco profundos, descalces, expectativas que se desanclan) y como complemento de la flotación, no como sustituto del ajuste \parencite{FMI2023IPFIntervencion}.
  \item Es la versión moderna del trilema con esquinas redondeadas: un país que flota puede usar otros instrumentos para que la flotación no le haga daño.
  \end{itemize}
\end{itemize}

% --------------------------------------------------------------------------
\subsection{Regulación: los controles de capitales}\label{sec:regulacion}

\begin{itemize}
\item La intervención actúa sobre el precio de las divisas; la regulación, sobre quién puede comprarlas y para qué.
\item \emph{Conclusión anticipada}: los controles de capitales reducen el bienestar de un país pequeño sin fricciones, pero pueden mejorarlo si corrigen una externalidad; funcionan a corto plazo y cambian la composición de los flujos más que su volumen, y no sustituyen a unas políticas macroeconómicas adecuadas.
\end{itemize}

\subsubsection{Definición y formas}\label{sec:controles-formas}

\begin{itemize}
\item Los \concepto{controles de capitales} son las restricciones que impone un Gobierno a las operaciones financieras con el exterior. Al cambiar los flujos de divisas, afectan al tipo de cambio.
\item El FMI distingue 2 tipos de medidas \parencite{FMI2012VisionInstitucional}:
  \begin{itemize}
  \item las \concepto{medidas de gestión de los flujos de capital}, que limitan esos flujos y suelen discriminar según la residencia;
  \item las \concepto{medidas macroprudenciales}, que limitan riesgos del sistema financiero (por ejemplo, el endeudamiento en divisas) sin mirar quién es residente \vertema{3.B.27}.
  \end{itemize}
\item La regulación de la conducta en el mercado de divisas (el Código Global de Conducta, la liquidación en CLS) es otra cosa: busca un mercado íntegro, no un tipo de cambio \vertema{3.B.13}.
\item \emph{Las formas}:
  \begin{itemize}
  \item \emph{cuantitativas}: prohibiciones, límites o autorizaciones previas, por ejemplo a la compra de inmuebles por no residentes, a la deuda exterior de los bancos o a la repatriación de beneficios; en las crisis, límites a la retirada de depósitos y a las transferencias al exterior;
  \item \emph{de precio}: impuestos sobre las operaciones o depósitos obligatorios sin remunerar, que encarecen sobre todo los flujos a corto plazo;
  \item \emph{tipos de cambio múltiples}: un tipo distinto según la operación.
  \end{itemize}
\item \emph{Los impuestos}:
  \begin{itemize}
  \item el \concepto{impuesto de \autor{Tobin}}: un pequeño impuesto sobre todas las operaciones de cambio, para «echar arena en los engranajes» de la especulación a corto plazo \parencite{Tobin1978Proposal}. Exige que lo apliquen todos los países, o las operaciones se van a otra plaza;
  \item el \concepto{encaje} chileno (1991-1998): un depósito sin remunerar en el banco central de una parte de cada entrada de capital (el 20~\% al principio, el 30~\% después) durante un año. Equivale a un impuesto tanto mayor cuanto más corto es el plazo de la entrada.
  \end{itemize}
\item \emph{Los tipos de cambio múltiples}: las autoridades fijan tipos distintos según la operación (las importaciones esenciales, las exportaciones, el capital), lo que exige un control de cambios que lleve cada operación a su mercado.
  \begin{itemize}
  \item El FMI los considera una práctica que exige su aprobación (artículo~VIII del Convenio Constitutivo), porque distorsionan los precios relativos y alimentan el mercado paralelo; el AREAER los recoge país por país \parencite{FMI2024AREAER2023}.
  \item Volvamos a España: tuvo tipos múltiples hasta el Plan de Estabilización de 1959, que los unificó en 60 pesetas por dólar \vertema{4.A.27}.
  \item Chile llegó a tener en 1973 15 tipos de cambio según la operación, con distorsiones que hundieron el comercio \parencite{Hudson1994Chile}; Venezuela vivió con un control de cambios y varios tipos oficiales (CADIVI, SICAD, DICOM) desde 2003, mientras el paralelo se disparaba.
  \end{itemize}
\end{itemize}

\subsubsection{Justificación y marco institucional}\label{sec:controles-justificacion}

\begin{itemize}
\item \emph{Sin fricciones}, la libertad de capitales es eficiente: lleva el ahorro adonde rinde más y permite suavizar el consumo (\ref{sec:controles-modelo}).
\item \emph{Con fricciones}, por la teoría del \emph{second best}, añadir una restricción puede mejorar el bienestar \vertema{3.A.22}. Las imperfecciones son condición necesaria, no suficiente:
  \begin{itemize}
  \item \emph{una externalidad pecuniaria}: quien se endeuda en divisas no tiene en cuenta que, si muchos lo hacen, una depreciación los obliga a vender a la vez y hunde los precios de las garantías de todos. Un impuesto a la entrada de deuda corrige el exceso, como un impuesto pigouviano \parencite{Korinek2011New} \vertema{3.B.13} \vertema{3.B.27};
  \item \emph{la sostenibilidad de la deuda}: abrir la cuenta financiera abarata el crédito y puede disparar el consumo y la deuda exterior, sobre todo con instituciones frágiles; los controles frenan los flujos y dan tiempo;
  \item \emph{el trilema}: un país con tipo fijo o intermedio que quiere una política monetaria propia necesita frenar las entradas (o las salidas) de capital, como en Bretton Woods (\ref{sec:trilema});
  \item \emph{un país grande}: si su demanda de fondos mueve el tipo de interés mundial, limitar sus préstamos lo baja y mejora su relación de intercambio intertemporal, como un arancel óptimo, si los demás no responden (anexo~\ref{anexo:controles}).
  \end{itemize}
\item \emph{El marco institucional}:
  \begin{itemize}
  \item el Convenio Constitutivo del FMI permite controlar los movimientos de capital (artículo~VI), pero no restringir los pagos corrientes (artículo~VIII) \vertema{3.B.22};
  \item la \concepto{visión institucional} del FMI, de 2012, acepta los controles en ciertas circunstancias: la liberalización debe ser gradual y ordenada; ante una oleada de entradas, primero las políticas macroeconómicas, y las medidas sobre los flujos solo como complemento, temporales y transparentes \parencite{FMI2012VisionInstitucional};
  \item su revisión de 2022 admite usarlas de forma preventiva sobre las entradas de deuda, aunque no haya oleada, si hay vulnerabilidades como un \emph{stock} grande de deuda en divisas \parencite{FMI2022RevisionVI};
  \item en la UE, la libre circulación de capitales es una libertad del Tratado (artículo~63 del TFUE), con salvaguardias temporales frente a terceros países \vertema{3.B.39}.
  \item la OCDE tiene desde 1961 su Código de Liberalización de los Movimientos de Capital, que compromete a sus miembros a liberalizar de forma progresiva.
  \end{itemize}
\end{itemize}

\subsubsection{Efectos:\cortetitulo el modelo intertemporal con controles}\label{sec:controles-modelo}

\paragraph{a) Idea}

\begin{itemize}
\item Para ver qué hace un control de capitales hace falta un modelo con bienestar. El enfoque intertemporal de la cuenta corriente lo tiene: los flujos de capital son la otra cara de las decisiones de consumo e inversión en el tiempo \vertema{3.B.12}.
\item Lo usamos con un impuesto sobre el endeudamiento exterior \parencite{SchmittGroheUribeWoodford2022International}.
\end{itemize}

\paragraph{b) Modelo}

\subparagraph{i) Supuestos}

\begin{itemize}
\item El modelo de 2 periodos del enfoque intertemporal \vertema{3.B.12}, sin inversión para simplificar: un país pequeño recibe las dotaciones $Y_t$ e $Y_{t+1}$ y un agente representativo maximiza $U=u(C_t)+\beta\,u(C_{t+1})$, con $0<\beta<1$ el factor de descuento.
\item El país puede prestar y pedir prestado al tipo de interés real mundial $r^*$, y quiere pedir prestado: $r^*$ es menor que su tipo de autarquía.
\item El Gobierno grava cada unidad de deuda exterior con un impuesto $\tau$, no prohibitivo, y devuelve lo recaudado a los hogares en bloque.
\end{itemize}

\subparagraph{ii) Desarrollo}

\begin{itemize}
\item \emph{En autarquía}, el país consume su dotación: el punto $A$ del gráfico~\ref{fig:controles-capital}.
\item \emph{Con libertad de capitales}, el país consume en $C$, donde la curva de indiferencia $U_1$ es tangente a la recta nacional, que pasa por $A$ con pendiente $-(1+r^*)$; en $t$ tiene un déficit por cuenta corriente, $C_t-Y_t>0$.
\item \emph{Con el impuesto}, los hogares se endeudan al tipo interno $r=r^*+\tau$, y la ecuación de \autor{Euler} pasa a ser
  \begin{equation}\label{eq:euler-controles}
  u'(C_t)=\beta\,(1+r^*+\tau)\,u'(C_{t+1}).
  \end{equation}
  \begin{itemize}
  \item Como el impuesto vuelve a los hogares, el país sigue pagando $r^*$ al exterior: el equilibrio está sobre la recta nacional, en el punto $D$, entre $A$ y $C$.
  \item En $D$, la curva de indiferencia $U_2$ tiene la pendiente del tipo interno, $-(1+r^*+\tau)$, más inclinada que la recta nacional (anexo~\ref{anexo:controles}).
  \end{itemize}
\end{itemize}

\begin{figure}[H]
\centering
\caption{El efecto de un impuesto sobre el endeudamiento exterior}\label{fig:controles-capital}
\grafico[0.62\textwidth]{controles-capital}
\fuente{Elaboración propia a partir de \textcite{SchmittGroheUribeWoodford2022International} y \textcite[cap.~1]{ObstfeldRogoff1996Foundations}. $A$ es la dotación (y la autarquía); $C$, el consumo con libertad de capitales; $D$, con el impuesto $\tau$.}
\end{figure}

\subparagraph{iii) Implicaciones}

\begin{itemize}
\item \emph{Sobre el consumo}: el impuesto encarece consumir hoy; el consumo presente baja y el futuro sube.
\item \emph{Sobre la inversión}: si la hubiera, también bajaría, porque se invierte hasta que el capital rinde $r^*+\tau$ en vez de $r^*$ \vertema{3.B.12}.
\item \emph{Sobre la cuenta corriente}: con menos consumo y menos inversión, el déficit corriente de $t$ se reduce. El control puede corregir un desequilibrio exterior.
\item \emph{Sobre el bienestar}: $U_0<U_2<U_1$. El país pide prestado menos de lo óptimo y está peor que con libertad de capitales, aunque mejor que en autarquía: en un país pequeño sin fricciones, el control es una distorsión, como un arancel.
\end{itemize}

\paragraph{c) Extensiones}

\begin{itemize}
\item \emph{Con una externalidad}: si endeudarse en divisas tiene una externalidad pecuniaria, el nivel privado de deuda es excesivo y un impuesto bien elegido mejora el bienestar \parencite{Korinek2011New}. Es lo que da la vuelta al resultado y justifica la visión institucional del FMI.
\item \emph{Un país grande}: además de la distorsión, el control baja el tipo de interés mundial y el coste de su deuda, y puede salir ganando, salvo que otro país grande responda y ambos acaben peor, como en una guerra arancelaria (anexo~\ref{anexo:controles}).
\end{itemize}

\begin{anotaciones}
\begin{itemize}
\item El modelo, con el manual y sus diapositivas: \parencite{SchmittGroheUribeWoodford2022International,Uribe2022DiapositivasControles}.
\end{itemize}
\end{anotaciones}

\paragraph{d) Valoración}

\begin{itemize}
\item El modelo explica por qué la libertad de capitales es la regla y los controles la excepción: sin fricciones, solo distorsionan.
\item Pero es un modelo sin fricciones financieras, y las razones para controlar los capitales están precisamente en ellas. Lo que queda es saber si los controles funcionan en la práctica.
\end{itemize}

\subsubsection{Experiencias y evidencia}\label{sec:controles-evidencia}

\begin{itemize}
\item \emph{Las experiencias}:
  \begin{itemize}
  \item \emph{Chile} (1991-1998): el encaje sobre las entradas de capital, el ejemplo de control a la entrada más estudiado;
  \item \emph{Malasia} (1998): controles a la salida y un tipo fijo de 3,80 ringgit por dólar en plena crisis asiática, a contracorriente de la opinión del FMI de entonces \parencite{Singh2023Malaysia};
  \item \emph{Brasil} (2009-2013): un impuesto a la entrada de inversión de cartera ante la fuerte apreciación del real \parencite{ForbesFratzscherStraub2015Capital};
  \item \emph{Islandia} (2008-2017), \emph{Chipre} (2013-2015) y \emph{Grecia} (2015-2019): controles a la salida para frenar la huida de depósitos en una crisis bancaria; en Chipre y Grecia, dentro del euro \parencite{FMI2022RevisionVI} \vertema{3.B.16};
  \item Volvamos a España: en septiembre de 1992, en plena crisis del SME, España, Irlanda y Portugal recurrieron a controles temporales para defender sus monedas \parencite{EichengreenWyplosz1993Unstable}.
  \end{itemize}
\item \emph{La evidencia}:
  \begin{itemize}
  \item los controles a la entrada dan algo más de autonomía monetaria, alargan los plazos de los flujos y reducen algo las presiones sobre el tipo real, pero no reducen el volumen neto de los flujos; los de salida de Malasia sí frenaron las salidas. La literatura se apoyaba demasiado en Chile y Malasia \parencite{MagudReinhartRogoff2018Capital};
  \item las medidas macroprudenciales reducen la fragilidad financiera, pero ni ellas ni las medidas sobre los flujos cambian mucho el tipo de cambio ni los flujos agregados \parencite{ForbesFratzscherStraub2015Capital}.
  \end{itemize}
\item \emph{El balance de los controles}:
  \begin{itemize}
  \item \emph{a favor}: reducen la volatilidad del tipo de cambio al penalizar la especulación a corto plazo, y dan tiempo para corregir los desequilibrios internos;
  \item \emph{en contra}: desincentivan la inversión, frenan también el arbitraje (y con él la integración de los mercados y la ley del precio único), distorsionan las decisiones intertemporales y se esquivan con el tiempo.
  \end{itemize}
\item Recapitulando el apartado: la intervención y los controles son los instrumentos con que se sostiene un régimen y se redondean las esquinas del trilema; funcionan como complemento de unas políticas coherentes con el régimen, nunca como sustituto.
\end{itemize}

% ==========================================================================
\section*{Conclusión}
\minutos{1}

\begin{itemize}
\item \textbf{\emph{Recapitulación}}:
  \begin{itemize}
  \item El trilema obliga a elegir: con libertad de capitales, fijar el tipo de cambio es renunciar a la política monetaria propia, y flotar no aísla del todo del ciclo financiero global.
  \item El tipo fijo da credibilidad y estabilidad al comercio, pero deja el ajuste a precios y salarios y expone a ataques; el flexible absorbe las perturbaciones reales, pero trae volatilidad y, con deuda en dólares, miedo a flotar.
  \item Entre los 2 polos hay un menú de regímenes, y la mitad de los países está en medio; sobreviven los intermedios que se defienden con credibilidad y reservas.
  \item Para sostener el régimen, los países intervienen y a veces limitan los capitales: los 2 instrumentos funcionan como complemento, no como sustituto, de las políticas adecuadas.
  \end{itemize}
\item \textbf{\emph{Relevancia}}:
  \begin{itemize}
  \item Para España, la elección se hizo con el euro: renunció al tipo de cambio dentro de la zona y comparte una moneda que flota frente al resto del mundo \vertema{3.B.16} \vertema{4.A.28}.
  \end{itemize}
\item \textbf{\emph{Extensiones}}:
  \begin{itemize}
  \item La determinación del tipo de cambio \vertema{3.B.14}, las crisis cambiarias que acaban con los tipos fijos \vertema{3.B.17}, la coordinación internacional y las guerras de divisas \vertema{3.B.19}, la historia del sistema monetario internacional \vertema{3.B.20} \vertema{3.B.21} y el FMI \vertema{3.B.22}.
  \end{itemize}
\item \textbf{\emph{Opinión}}:
  \begin{itemize}
  \item La paradoja de los últimos 20 años es que, mientras los tipos fijos perdían peso en la economía mundial, el consenso aceptaba la intervención y los controles como complemento de la flotación \parencite{FMI2020IPF,FMI2022RevisionVI}. No es una vuelta atrás: es reconocer que el trilema tiene las esquinas redondeadas.
  \end{itemize}
\item \textbf{\emph{Idea final}}:
  \begin{itemize}
  \item \textcite{Frankel1999NoSingle} lo resumió en el título de su ensayo: ningún régimen sirve a todos los países ni en todo momento. Cada uno sabe dónde le aprieta el zapato.
  \item Lo que ningún país puede hacer es elegir un régimen y aplicar políticas que no casen con él: eso es el trilema, y todo lo demás es comentario.
  \end{itemize}
\end{itemize}

\clearpage
\appendix

% ==========================================================================
\section{Anexo: los regímenes del FMI, uno a uno}\label{anexo:regimenes}

\begin{itemize}
\item Este anexo recorre los 10 regímenes de la clasificación de hecho del FMI con su definición, sus ventajas e inconvenientes y un ejemplo vigente, con su serie cuando la hay \parencite{HabermeierEtAl2009Revised,FMI2026InformeAnual}.
\item Número de economías en cada régimen a 30 de abril de 2026 (clasificación preliminar, sobre 195): sin moneda propia, 15; caja de conversión, 12; convencional, 40; bandas horizontales, 1; estabilizado, 25; deslizamiento, 1; similar al deslizamiento, 25; flotación, 34; flotación libre, 31, y otros gestionados, 11 \parencite{FMI2026InformeAnual}.
\item Los miembros de una unión monetaria se clasifican por el régimen de la moneda común: la zona del euro, en libre flotación; la UEMOA y la CEMAC, en el anclaje convencional, y el Caribe Oriental, en caja de conversión. La dolarización oficial, unilateral o con acuerdo, es la categoría «sin moneda propia de curso legal»; la no oficial no es un régimen.
\end{itemize}

\subsection{Anclajes duros}\label{anexo:regimenes-duros}

\begin{itemize}
\item \emph{Sin moneda propia de curso legal}: Ecuador, El Salvador y Panamá con el dólar; Kosovo, Montenegro, San Marino y Andorra con el euro. Ventajas e inconvenientes en~\ref{sec:anclajes-duros}.
\item \emph{Caja de conversión}: el dólar de Hong Kong se mueve dentro de la banda de 7,75 a 7,85 por dólar desde 2005, y ha pasado años cerca de cada uno de los bordes (gráfico~\ref{fig:hong-kong}). Bulgaria, en el euro desde el 1 de enero de 2026, todavía figura como caja de conversión en la clasificación preliminar de abril de 2026 \parencite{FMI2026InformeAnual}.
\end{itemize}

\begin{figure}[H]
\centering
\caption{Caja de conversión: el dólar de Hong Kong, 2000-2026 (HKD por USD)}\label{fig:hong-kong}
\grafico[0.66\textwidth]{hong-kong-caja-conversion}
\fuente{Elaboración propia con datos de la Reserva Federal \parencite{Fed2026H10}. Media mensual; la banda de 7,75 a 7,85, desde mayo de 2005 \parencite{HKMA2005Refinamientos}.}
\end{figure}

\subsection{Anclajes blandos}\label{anexo:regimenes-blandos}

\begin{itemize}
\item \emph{Convencional}: el país ancla su moneda, \emph{de jure} o \emph{de facto}, a otra o a una cesta (por ejemplo, la de sus principales socios, para fijar el tipo efectivo) y la deja fluctuar menos del $\pm1$~\%; notifica el ancla al FMI y defiende la paridad con intervención directa.
  \begin{itemize}
  \item \emph{Ventajas}: credibilidad y estabilidad para el comercio con el ancla, sin renunciar a la moneda.
  \item \emph{Inconvenientes}: sin política monetaria propia si el capital es libre, y expuesto a ataques si la paridad deja de ser creíble.
  \item \emph{Ejemplo}: el riyal saudí, a 3,75 por dólar desde 1986 (gráfico~\ref{fig:riyal}).
  \end{itemize}
\end{itemize}

\begin{figure}[H]
\centering
\caption{Paridad fija: el riyal de Arabia Saudí, 2000-2026 (SAR por USD)}\label{fig:riyal}
\grafico[0.66\textwidth]{riyal-paridad-fija}
\fuente{Elaboración propia con datos del Banco de Pagos Internacionales \parencite{BIS2026TiposBilaterales}. Media mensual.}
\end{figure}

\begin{itemize}
\item \emph{Estabilizado}: el tipo se mantiene en un margen del 2~\% durante al menos 6 meses por acción oficial (incluidas las rigideces del mercado), sin compromiso político de las autoridades.
  \begin{itemize}
  \item \emph{Ventajas}: la estabilidad de un tipo fijo sin atarse públicamente a él.
  \item \emph{Inconvenientes}: sin el compromiso, no importa toda la credibilidad del ancla, y la estabilidad puede romperse cuando deja de interesar.
  \item \emph{Ejemplos}: el denar de Macedonia del Norte, en torno a 61,5 por euro desde 2002 con un régimen oficial de flotación (gráfico~\ref{fig:denar}); Bolivia, con un deslizamiento oficial que en la práctica es un ancla al dólar, y Nicaragua desde 2024. La kuna croata, que estaba en este grupo, desapareció con la entrada de Croacia en el euro el 1 de enero de 2023 \parencite{FMI2024AREAER2023}.
  \end{itemize}
\end{itemize}

\begin{figure}[H]
\centering
\caption{Régimen estabilizado: el denar de Macedonia del Norte, 2002-2026 (MKD por EUR)}\label{fig:denar}
\grafico[0.66\textwidth]{denar-estabilizado}
\fuente{Elaboración propia con datos del Banco de Pagos Internacionales \parencite{BIS2026TiposBilaterales}. Media mensual, cociente de los tipos del denar y del euro frente al dólar.}
\end{figure}

\begin{itemize}
\item \emph{Deslizamiento}: la paridad se ajusta en pequeñas cantidades, a un ritmo anunciado de antemano o según indicadores como la inflación pasada; puede ser de depreciación o de apreciación, y sus variantes difieren en la publicidad de la senda, su velocidad y la frecuencia con que se revisa.
  \begin{itemize}
  \item \emph{Ventajas}: más flexibilidad que un tipo fijo, porque admite un diferencial de inflación con el ancla sin que el tipo real se aprecie, y algo más de autonomía monetaria.
  \item \emph{Inconvenientes}: el margen de política monetaria es muy reducido, y la apreciación real no desaparece si la inflación supera al deslizamiento.
  \item \emph{Ejemplo}: el córdoba de Nicaragua, que se depreció a un ritmo anunciado desde 1993 hasta que el deslizamiento se puso a cero el 1 de enero de 2024 (gráfico~\ref{fig:cordoba}) \parencite{BCN2023Deslizamiento}. Hoy el único deslizamiento es el de Botsuana, con una cesta de rand y derechos especiales de giro \parencite{FMI2026InformeAnual}.
  \end{itemize}
\end{itemize}

\begin{figure}[H]
\centering
\caption{Deslizamiento: el córdoba de Nicaragua, 2000-2026 (NIO por USD)}\label{fig:cordoba}
\grafico[0.66\textwidth]{cordoba-deslizamiento}
\fuente{Elaboración propia con datos del Banco de Pagos Internacionales \parencite{BIS2026TiposBilaterales} y del Banco Central de Nicaragua \parencite{BCN2023Deslizamiento}. Media mensual; desde enero de 2024, 36,6243.}
\end{figure}

\begin{itemize}
\item \emph{Similar al deslizamiento}: el tipo se mantiene en un margen del 2~\% en torno a una tendencia identificada estadísticamente durante al menos 6 meses, sin que pueda considerarse flotación.
  \begin{itemize}
  \item \emph{Ventajas e inconvenientes}: los del deslizamiento, sin el compromiso de anunciar la senda.
  \item \emph{Ejemplo}: el renminbi chino, que figura en este grupo según la clasificación preliminar de 2026, tras pasar por «otros gestionados» desde marzo de 2022 (gráfico~\ref{fig:renminbi}) \parencite{FMI2024AREAER2023,FMI2026InformeAnual}.
  \end{itemize}
\end{itemize}

\begin{figure}[H]
\centering
\caption{Régimen similar al deslizamiento: el renminbi chino, 2000-2026 (CNY por USD)}\label{fig:renminbi}
\grafico[0.66\textwidth]{renminbi-crawl-like}
\fuente{Elaboración propia con datos de la Reserva Federal \parencite{Fed2026H10}. Media mensual.}
\end{figure}

\begin{itemize}
\item \emph{Bandas horizontales}: el tipo se mueve en una banda de al menos el $\pm1$~\% en torno a una paridad central, o entre un máximo y un mínimo que difieren más del 2~\%. Es el régimen de Marruecos y el de las zonas objetivo (\ref{sec:krugman}).
  \begin{itemize}
  \item \emph{Ventajas}: si es creíble, ancla las expectativas y reduce la volatilidad, y deja algo de autonomía monetaria.
  \item \emph{Inconvenientes}: los realineamientos aumentan la volatilidad, y una banda que no se cree invita al ataque.
  \end{itemize}
\end{itemize}

\subsection{Flotación y otros regímenes}\label{anexo:regimenes-flotacion}

\begin{itemize}
\item \emph{Flotación}: el dólar neozelandés y el real brasileño flotan con intervenciones ocasionales (gráfico~\ref{fig:flotacion}); Egipto, desde marzo de 2024, cuando dejó flotar la libra tras unificar sus tipos (gráfico~\ref{fig:libra-egipcia}) \parencite{CBE2024Flotacion}.
\item \emph{Flotación libre}: Estados Unidos, la zona del euro, Japón, el Reino Unido, Canadá, Australia o México \parencite{FMI2026InformeAnual}.
\item \emph{Otros gestionados}: la categoría residual, para los regímenes que no encajan en ninguna otra o que cambian con frecuencia: Argentina, Irán, Kuwait, Etiopía o Venezuela \parencite{FMI2026InformeAnual}.
\end{itemize}

\begin{figure}[H]
\centering
\caption{Flotación: el dólar neozelandés y el real brasileño, 2000-2026 (por USD)}\label{fig:flotacion}
\grafico[0.8\textwidth]{flotacion-nzd-brl}
\fuente{Elaboración propia con datos de la Reserva Federal \parencite{Fed2026H10}. Medias mensuales.}
\end{figure}

\begin{figure}[H]
\centering
\caption{De un anclaje a la flotación: la libra egipcia, 2015-2026 (EGP por USD)}\label{fig:libra-egipcia}
\grafico[0.66\textwidth]{libra-egipcia-flotacion}
\fuente{Elaboración propia con datos del Banco de Pagos Internacionales \parencite{BIS2026TiposBilaterales}. Media mensual; la línea vertical marca la unificación del 6 de marzo de 2024 \parencite{CBE2024Flotacion}.}
\end{figure}

\begin{anotaciones}
\begin{itemize}
\item El vídeo del FMI que explica la clasificación con estos ejemplos: \parencite{FMIInstitute2021Clasificacion}.
\item Cómo clasificaba el FMI los regímenes en 2004, con el paso a la clasificación de hecho: \parencite{FMI2004Clasificacion}.
\end{itemize}
\end{anotaciones}

% ==========================================================================
\section{Anexo: las zonas objetivo de \autor{Krugman}, paso a paso}\label{anexo:krugman}

\begin{itemize}
\item \emph{El punto de partida} es~\eqref{eq:zona-objetivo}, con los fundamentos $f_t\equiv m_t+v_t$. Dentro de la banda, $m_t$ no cambia y $\mathrm{d}v_t=\sigma\,\mathrm{d}z_t$, un paseo aleatorio en tiempo continuo sin deriva ($z_t$ es un movimiento browniano).
\item \emph{Una solución que depende solo de los fundamentos}: buscamos $e_t=g(f_t)$. Por el lema de \autor{Itô}, $\mathbb{E}_t[\mathrm{d}e_t]/\mathrm{d}t=\tfrac{\sigma^2}{2}\,g''(f_t)$, y~\eqref{eq:zona-objetivo} queda
  $$g(f)=f+\frac{\alpha\sigma^2}{2}\,g''(f).$$
\item \emph{La solución general} es la recta de 45° más 2 exponenciales:
  $$g(f)=f+A\,\mathrm{e}^{\lambda f}+B\,\mathrm{e}^{-\lambda f},\qquad \lambda=\sqrt{\frac{2}{\alpha\sigma^2}}.$$
\item \emph{Una banda simétrica} en torno a la paridad central (que normalizamos en $f=0$) exige $B=-A$:
  $$g(f)=f+A\left(\mathrm{e}^{\lambda f}-\mathrm{e}^{-\lambda f}\right).$$
\item \emph{El \emph{smooth pasting}}: en los bordes de los fundamentos, $\pm\bar f$, donde la autoridad interviene, la curva tiene que ser tangente a la banda, $g'(\bar f)=0$, porque si no los agentes tendrían una ganancia segura al acercarse al borde:
  $$1+A\lambda\left(\mathrm{e}^{\lambda\bar f}+\mathrm{e}^{-\lambda\bar f}\right)=0\quad\Rightarrow\quad A=-\frac{1}{\lambda\left(\mathrm{e}^{\lambda\bar f}+\mathrm{e}^{-\lambda\bar f}\right)}<0.$$
\item \emph{Por qué es una S aplanada}: con $A<0$, para $f>0$ el término exponencial es negativo y $g(f)<f$; para $f<0$, al revés. El tipo queda dentro de la recta de 45°, y los bordes de la banda son $e^{\max}=g(\bar f)$ y $e^{\min}=-g(\bar f)$ \parencite{Krugman1991Target}. \findemo
\item \emph{Sin credibilidad}: si en el borde hay una probabilidad de realineamiento en vez de defensa, el signo de $A$ puede cambiar y la curva se vuelve una \emph{S} invertida, más inclinada que la de 45° \parencite{BertolaCaballero1992Target}.
\end{itemize}

% ==========================================================================
\section{Anexo: el control de capitales, paso a paso}\label{anexo:controles}

\begin{itemize}
\item \emph{Los hogares}: piden prestado $B_t=C_t-Y_t>0$ en $t$ y devuelven $(1+r^*+\tau)B_t$ en $t+1$; reciben la transferencia $T_{t+1}$, que toman como dada. Su problema es
  $$\optim{\max}{C_t,\,C_{t+1}}{u(C_t)+\beta\,u(C_{t+1})}{C_{t+1}=Y_{t+1}-(1+r^*+\tau)\,(C_t-Y_t)+T_{t+1}}$$
\item \emph{La condición}: sustituyendo la restricción en el objetivo y derivando respecto a $C_t$, $u'(C_t)-\beta\,(1+r^*+\tau)\,u'(C_{t+1})=0$, que es~\eqref{eq:euler-controles}.
\item \emph{El equilibrio}: el Gobierno devuelve lo que recauda, $T_{t+1}=\tau B_t$. Sustituyendo en la restricción, $C_{t+1}=Y_{t+1}-(1+r^*)\,(C_t-Y_t)$: el país en su conjunto paga $r^*$ al exterior, y su restricción es la nacional,
  $$C_t+\frac{C_{t+1}}{1+r^*}=Y_t+\frac{Y_{t+1}}{1+r^*}.$$
\item \emph{El punto $D$}: está sobre la recta nacional, donde la relación marginal de sustitución es $1+r^*+\tau$ en vez de $1+r^*$. Como la curva de indiferencia es convexa, $D$ queda entre $A$ y $C$: menos endeudamiento y menos déficit corriente que en $C$. \findemo
\item \emph{El bienestar}: $C$ es el mejor punto de la recta nacional, así que $U_2<U_1$. Y como la utilidad crece a lo largo de la recta de $A$ a $C$ (las preferencias son convexas), $D$, que está entre los 2, da más utilidad que $A$: $U_2>U_0$.
\item \emph{Un país grande}: si su endeudamiento sube el tipo mundial, $r^*=r^*(B_t)$ con $r^{*\prime}>0$, cada unidad adicional de deuda encarece toda la anterior. Los hogares no lo tienen en cuenta, y un impuesto $\tau=B_t\,r^{*\prime}(B_t)$ les hace internalizar ese coste, como un arancel óptimo.
  \begin{itemize}
  \item El país puede estar mejor con el control que con libertad de capitales, según los parámetros, aunque haya represalias.
  \item Pero si un país grande impone su control óptimo, el otro tiene interés en responder, y los 2 pueden acabar peor que sin controles, como en una guerra arancelaria.
  \end{itemize}
\end{itemize}

% ==========================================================================
\section{Anexo: para el test}\label{anexo:test}

\begin{itemize}
\item \emph{Fechas y obras}: \textcite{Friedman1953Case}, a favor de los tipos flexibles; \textcite{Mundell1961Theory}, áreas monetarias óptimas; \textcite{Fleming1962Domestic} y \textcite{Mundell1963Capital}, el modelo que da el trilema; \textcite{Krugman1991Target}, zonas objetivo; \textcite{CalvoReinhart2002}, miedo a flotar; \textcite{Rey2013Dilema}, el dilema; Acuerdo Smithsonian (diciembre de 1971), bandas del $\pm2{,}25$~\%; flotación generalizada desde marzo de 1973; Segunda Enmienda del Convenio del FMI, en vigor desde 1978; visión institucional del FMI (2012, revisada en 2022); Marco de Política Integrado (2020).
\item \emph{Trampas}:
  \begin{itemize}
  \item la BP vertical es movilidad \emph{nula} del capital ($\beta_K=0$), no perfecta;
  \item la zona del euro está en \emph{libre flotación} en la clasificación del FMI, y la unión monetaria no es una categoría aparte;
  \item Montenegro y Kosovo no tienen moneda propia: usan el euro como única moneda de curso legal (dolarización oficial y unilateral);
  \item en la flotación libre se interviene como mucho 3 veces en 6 meses, de no más de 3 días hábiles cada una;
  \item esterilizar una compra de divisas es \emph{vender} títulos;
  \item la intervención esterilizada funciona por el canal de cartera si los activos no son sustitutos perfectos, y por el de señales aunque lo sean.
  \end{itemize}
\end{itemize}

\begin{anotaciones}
Preguntas de test de exámenes anteriores, con su explicación: \url{https://www.victorgutierrezmarcos.es/oposicion/temario/primer-ejercicio/test/simulador.html}.
\end{anotaciones}

\bibliografia

\end{document}
